\documentclass[10pt,a4paper]{article}
\usepackage[utf8]{inputenc}
\usepackage[T1]{fontenc}
\usepackage{amsmath,amssymb,amsfonts}
\usepackage{geometry}
\geometry{margin=2.5cm}
\usepackage{booktabs}
\usepackage{array}
\usepackage{longtable}
\usepackage{xcolor}
\usepackage{hyperref}
\hypersetup{colorlinks=true,linkcolor=blue,urlcolor=blue}
\usepackage{enumitem}
\usepackage{graphicx}
\usepackage{float}
\usepackage{verbatim}
\usepackage{tabularx}
% \usepackage{microtype}
\usepackage{pifont}
\usepackage{appendix}
\usepackage{textcomp}
\usepackage{pdflscape}
\newcommand{\cmark}{\ding{51}}
\setlength{\parindent}{0pt}
\setlength{\parskip}{6pt}

\title{\textbf{Wing Structure Design Calculations}\\[0.5em]
\Large Solar UAV Wing \\[10pt]
Group~2, IIST}
\author{}
\date{}

\begin{document}
\maketitle
\tableofcontents
\newpage

\newpage \section{Introduction and Use Case}

The aircraft is a \textbf{solar-powered UAV} designed for \textbf{continuous, long-endurance flight} (24 h+), scaled from the Sky-Sailor design methodology (Noth \& Siegwart, ETH 2008). The mission profile drives every structural choice:

\begin{center}
\begin{tabular}{|c|c|}
\hline
\textbf{Mission requirement} & \textbf{Structural implication} \\
\hline
Continuous solar flight & Ultra-low structural mass \\
\hline
Solar-cell upper surface & Flat, non-buckling upper skin ${\rightarrow}$ many ribs \\
\hline
Low wing loading & High span, thin airfoil ${\rightarrow}$ bending-dominated spar \\
\hline
Low Reynolds number (Re ${\approx}$ 1.0 ${\times}$ 10$^{5}$) & Careful surface quality; airfoil shape held by ribs \\
\hline
Hand-launched, belly-land & Root loads govern sizing; no landing gear loads \\
\hline
\end{tabular}
\end{center}

\textbf{Reference aircraft - Sky-Sailor (Noth 2008):}

\begin{center}
\begin{tabular}{|l|l|}
\hline
Parameter & Value \\
\hline
Wingspan & 3.20 m \\
\hline
Wing area & 0.776 m$^{2}$ \\
\hline
Aspect ratio & 13.2 \\
\hline
Total mass & 2.44 kg \\
\hline
Airframe mass & 0.725 kg \\
\hline
Cruise speed & 8.3 m/s \\
\hline
Cruise Re & ~150000 \\
\hline
\end{tabular}
\end{center}

Our design is a \textbf{geometrically scaled-down} laboratory wing/prototype using the same WE3.55-9.3 airfoil, with a span of 980 mm.




\newpage \section{Wing Geometry and Airfoil Parameters}

\subsection{Wing planform}


\begin{center}
\begin{tabular}{|l|l|l|}
\hline
Parameter & Symbol & Value \\
\hline
Full span & b & 980 mm = 0.980 m \\
\hline
Half-span & b/2 & 490 mm = 0.490 m \\
\hline
Chord (constant) & c & 250 mm = 0.250 m \\
\hline
Wing planform area & S & b ${\times}$ c = 0.245 m$^{2}$ \\
\hline
Aspect ratio & AR & b$^{2}$/S = 0.980$^{2}$/0.245 = \textbf{3.92} \\
\hline
Taper ratio & ${\lambda}$ & 1.00 (rectangular planform) \\
\hline
Dihedral & ${\Gamma}$ & 3$^{\circ}$ (assumed, for lateral stability) \\
\hline
Sweep & ${\Lambda}$ & 0$^{\circ}$ (unswept) \\
\hline
\end{tabular}
\end{center}

Wing area:
\[
S = b \times c = 0.980 \times 0.250 = 0.245\ \text{m}^2
\]

Aspect ratio:
\[
AR = \frac{b^2}{S} = \frac{(0.980)^2}{0.245} = \frac{0.9604}{0.245} = 3.92
\]

\begin{figure}[H]
    \centering
    \includegraphics[width=0.905\linewidth]{Fig1_Wing_Geometry_Schematic.png}
    \caption{Wing planform (top view). b = 980 mm, c = 250 mm and AR = 3.92.}
    \label{fig:fig1}
\end{figure}


\subsection{Airfoil: WE3.55-9.3}


\begin{center}
\begin{tabular}{|l|l|}
\hline
Parameter & Value \\
\hline
Max thickness & t/c = 9.3 \% \\
\hline
Max thickness location & ${\approx}$ 30 \% chord \\
\hline
Max camber & f/c = 3.55 \% \\
\hline
Max camber location & ${\approx}$ 40 \% chord \\
\hline
Design C\_L (cruise) & C\_L ${\approx}$ 0.80 \\
\hline
Pitching-moment coeff. (about AC) & C\_M,AC ${\approx}$ -0.05 \\
\hline
Aerodynamic centre & x\_AC/c ${\approx}$ 0.25 \\
\hline
Zero-lift angle of attack & ${\alpha}$\_0 ${\approx}$ -4$^{\circ}$ \\
\hline
\end{tabular}
\end{center}

\textbf{Structural depths at key chord-wise stations:}

\[
t_{\max} = \frac{t}{c} \times c = 0.093 \times 250\ \text{mm} = 23.25\ \text{mm}
\]

\begin{center}
\begin{tabular}{|l|l|l|l|}
\hline
Station & \% chord & x from LE (mm) & Airfoil thickness (approx, mm) \\
\hline
Leading edge & 0 \% & 0 & 0 \\
\hline
Front spar (FS) & 25 \% & 62.5 & 20.0 \\
\hline
Max thickness & 30 \% & 75.0 & 23.25 \\
\hline
Rear spar (RS) & 65 \% & 162.5 & 13.5 \\
\hline
Trailing edge & 100 \% & 250.0 & 0 \\
\hline
\end{tabular}
\end{center}

Thickness at the front spar (30\% of max thickness reached by 25\% chord, using NACA-like distribution; value interpolated):

\[
h_{FS} \approx 0.86 \times t_{\max} = 0.86 \times 23.25 = 20.0\ \text{mm}
\]

Thickness at the rear spar (approximately 58\% of max):

\[
h_{RS} \approx 0.58 \times t_{\max} = 0.58 \times 23.25 = 13.5\ \text{mm}
\]


\subsection{Spar and rib layout}


\begin{center}
\begin{tabular}{|l|l|l|l|}
\hline
Sub-structure & Location & Count & Material \\
\hline
Front (main) spar & 25 \% c = 62.5 mm from LE & 1 & Aluminium rod (solid circular) \\
\hline
Rear spar & 65 \% c = 162.5 mm from LE & 1 & Aluminium rod (solid circular) \\
\hline
Ribs & Spanwise, ~25 mm pitch & ${\approx}$ 40 & Balsa wood sheet \\
\hline
Leading-edge D-box skin & LE to FS & - & Aluminium sheet \\
\hline
Upper/lower skin & FS to RS (inter-spar bay) & - & Aluminium sheet (or film) \\
\hline
\end{tabular}
\end{center}






\newpage \section{Mass Estimation and Assumptions}

\subsection{Assumptions}

\begin{center}
\renewcommand{\arraystretch}{1.2}
\begin{tabularx}{\textwidth}{|p{4cm}|X|}
\hline
Assumption & Justification / Correction \\
\hline
\textbf{A1.} Total aircraft mass $m = 0.50$ kg & Scaled from Sky-Sailor using cubic scaling law (refer No.3.2). Conservative estimate for a 980 mm solar UAV. \\
\hline
\textbf{A2.} Design load factor $n = 2.5$ & Standard value for small UAVs (CS-23 small aircraft; acceptable for academic model). Gust correction applied in No.4.3. \\
\hline
\textbf{A3.} Ultimate factor of safety $\mathrm{FOS} = 1.5$ & FAR/CS standard; ensures plastic reserve beyond limit loads. \\
\hline
\textbf{A4.} Lift acts at $x_{CP} = 31.3\%\,c$ & Derived from $C_L$ and $C_M$ (refer No.4.2); corrected for off-design conditions. \\
\hline
\textbf{A5.} Rectangular wing $\rightarrow$ uniform load baseline & Elliptical correction applied in No.5.2. \\
\hline
\textbf{A6.} Cruise altitude $\approx$ sea level; $\rho = 1.225$ kg/m$^{3}$ & Low-altitude solar UAV; use ISA sea-level values. \\
\hline
\textbf{A7.} Aluminium alloy: Al 6061-T6 & Commonly available, good strength-to-weight; rod and sheet forms used. \\
\hline
\textbf{A8.} Balsa rib grain parallel to chord & Maximises bending stiffness in ribs; standard orientation. \\
\hline
\textbf{A9.} Single-cell D-box torsion model & Leading edge + front spar + upper/lower skin form closed D-box; rear-spar bay treated as open. \\
\hline
\textbf{A10.} Fixed-free cantilever beam model & Wing is cantilevered at root (centre fuselage joint); tip is free. \\
\hline
\end{tabularx}
\end{center}

\subsection{Total aircraft mass - scaling from Sky-Sailor}


Cubic scaling law (Noth Ch. 3, geometric similarity):

\[
m_{\text{our}} = m_{\text{SS}} \left(\frac{b_{\text{our}}}{b_{\text{SS}}}\right)^3 = 2.44 \times \left(\frac{0.980}{3.200}\right)^3
\]

\[
= 2.44 \times (0.3063)^3 = 2.44 \times 0.02873 = 0.0701\ \text{kg}
\]

This gives 70 g, which is unrealistically small for a functional prototype with aluminium components. The cubic scaling assumes geometric and material similarity (all-composite), which does not hold when switching to aluminium spars.

\textbf{Corrected estimate - structural weight model (Noth top-5\% model):}

\[
W_{af} = 0.44\, S^{1.55}\, AR^{1.30}\ [\text{N}]
\]

\[
W_{af} = 0.44 \times (0.245)^{1.55} \times (3.92)^{1.30}
\]

\[
W_{af} = 0.44 \times 0.1130 \times 5.907 = 0.294\ \text{N} \implies m_{af,\text{best}} = 30\ \text{g}
\]

This is the all-composite best-case prediction. For aluminium spars and balsa ribs, use the mean-quality model:

\[
W_{af,\text{mean}} = 5.58\, S^{1.59}\, AR^{0.71} = 5.58 \times (0.245)^{1.59} \times (3.92)^{0.71}
\]

\[
(0.245)^{1.59} = e^{1.59 \times (-1.407)} = e^{-2.237} = 0.1065
\]

\[
(3.92)^{0.71} = e^{0.71 \times 1.366} = e^{0.970} = 2.638
\]

\[
W_{af,\text{mean}} = 5.58 \times 0.1065 \times 2.638 = 1.567\ \text{N} \implies m_{af} \approx 160\ \text{g}
\]

\textbf{Adopted total aircraft mass}: Adding propulsion (~150 g), avionics (~100 g), battery (~100 g), structure (~160 g):

\[
\boxed{m_{\text{total}} = 0.50\ \text{kg}, \quad W = mg = 0.50 \times 9.81 = 4.905\ \text{N}}
\]

\textit{This is the design point. All structural calculations use this value.}

\newpage

\section{Aerodynamic Load Analysis}

\subsection{Cruise flight conditions}

Level cruise (load factor n = 1):

\[
L = W = 4.905\ \text{N}
\]

\[
\frac{W}{S} = \frac{4.905}{0.245} = 20.0\ \text{N/m}^2
\]

Cruise speed from the level-flight lift equation:

\[
v_{\text{cruise}} = \sqrt{\frac{2L}{\rho\, C_L\, S}} = \sqrt{\frac{2 \times 4.905}{1.225 \times 0.80 \times 0.245}}
\]

\[
= \sqrt{\frac{9.810}{0.2401}} = \sqrt{40.86} = 6.39\ \text{m/s}
\]

Reynolds number at cruise:

\[
Re = \frac{\rho\, v\, c}{\mu} = \frac{1.225 \times 6.39 \times 0.250}{1.789 \times 10^{-5}} = \frac{1.957}{1.789 \times 10^{-5}} = 1.09 \times 10^5
\]

This is in the valid operating range for WE3.55-9.3 (Re = 80000 - 300000).

Dynamic pressure at cruise:

\[
q_{\infty} = \frac{1}{2}\rho v^2 = \frac{1}{2} \times 1.225 \times (6.39)^2 = 0.6125 \times 40.83 = 25.01\ \text{N/m}^2
\]


\subsection{Centre of pressure location}


The aerodynamic centre is at x\_AC = 0.25c. The pitching moment about the aerodynamic centre is M\_AC = C\_M,AC ${\times}$ q ${\times}$ S ${\times}$ c. The lift acts at the centre of pressure:

\[
\frac{x_{CP}}{c} = \frac{x_{AC}}{c} - \frac{C_{M,AC}}{C_L} = 0.25 - \frac{-0.05}{0.80} = 0.25 + 0.0625 = 0.3125
\]

\[
x_{CP} = 0.3125 \times 250 = 78.1\ \text{mm from LE}
\]

\textbf{Eccentricity of lift from front spar:}

\[
e_L = x_{CP} - x_{FS} = 78.1 - 62.5 = 15.6\ \text{mm} = 0.0156\ \text{m}
\]


\subsection{Design loads (limit and ultimate)}


\textbf{Limit load factor} n = 2.5 (manoeuvre). This is the maximum steady pull-up expected.

\textbf{Gust load correction} (CS-23 simplified gust formula):

\[
\Delta n_{\text{gust}} = \frac{\rho\, U_e\, v\, a}{2 (W/S)}
\]

where U\_e = 9.1 m/s (gust speed at cruise; CS-23 for low-altitude), a = dC\_L/d${\alpha}$ per rad.

For finite wing: $a = \frac{2\pi AR}{AR + 2} = \frac{2\pi \times 3.92}{3.92 + 2} = \frac{24.63}{5.92} = 4.16\ \text{rad}^{-1}$

\[
\Delta n_{\text{gust}} = \frac{1.225 \times 9.1 \times 6.39 \times 4.16}{2 \times 20.0} = \frac{292.3}{40.0} = 7.31
\]

\textbf{Gust alleviation factor (FAR 23 No.23.341):} The raw gust increment overestimates the effective load because inertia of the aircraft partially attenuates the instantaneous gust response. The alleviation factor K\_g is:

\[
K_g = \frac{0.88\,\mu_g}{5.3 + \mu_g}
\]

where the mass ratio $\mu_g$ is:

\[
\mu_g = \frac{2\,(m/S)}{\rho\,c\,a} = \frac{2 \times (0.500/0.245)}{1.225 \times 0.250 \times 4.16} = \frac{2 \times 2.041}{1.274} = \frac{4.082}{1.274} = 3.204
\]

\[
K_g = \frac{0.88 \times 3.204}{5.3 + 3.204} = \frac{2.819}{8.504} = 0.332
\]

Alleviated gust increment:

\[
\Delta n_{\text{gust,allev}} = K_g \times \Delta n_{\text{gust}} = 0.332 \times 7.31 = 2.43
\]

\[
n_{\text{gust,total}} = 1 + 2.43 = 3.43
\]

\textbf{Conclusion:} After gust alleviation, $n_{\text{gust}} = 3.43 \approx n_{\text{manoeuvre}} = 2.5$; the manoeuvre limit governs. The structural design uses \textbf{n = 2.5} as the limit load factor.

Even the alleviated gust load factor (3.43) exceeds the 2.5 manoeuvre limit. This is typical for ultra-light solar UAVs with very low wing loading (W/S = 20 N/m$^{2}$). Operationally, flights must be restricted to calm conditions (winds ${\leq}$ 5 m/s) to keep the effective load factor within the structural envelope.

\textbf{Limit lift force (full wing):}

\[
L_{\text{limit}} = n \times W = 2.5 \times 4.905 = 12.26\ \text{N}
\]

\textbf{Ultimate lift force:}

\[
L_{\text{ult}} = FOS \times L_{\text{limit}} = 1.5 \times 12.26 = 18.39\ \text{N}
\]

\textbf{Per half-wing (single panel, cantilevered at root):}

\[
L_{1/2,\text{limit}} = \frac{L_{\text{limit}}}{2} = 6.13\ \text{N}
\]

\[
L_{1/2,\text{ult}} = \frac{L_{\text{ult}}}{2} = 9.20\ \text{N}
\]


\subsection{Flight Envelope (V-n Diagram)}


The V-n diagram defines the structural flight envelope - the combinations of speed and load factor the wing is designed to withstand.

\textbf{Stall speed:}

\[
V_{\text{stall}} = \sqrt{\frac{2W}{\rho\,C_{L,\max}\,S}} = \sqrt{\frac{2 \times 4.905}{1.225 \times 1.25 \times 0.245}} = \sqrt{\frac{9.810}{0.3756}} = \sqrt{26.12} = 5.11\ \text{m/s}
\]

where $C_{L,\max} = 1.25$ estimated for the WE3.55-9.3 airfoil at Re ${\approx}$ 100,000.

\textbf{Manoeuvre speed:}

\[
V_A = V_{\text{stall}} \times \sqrt{n_{\text{max}}} = 5.11 \times \sqrt{2.5} = 5.11 \times 1.581 = 8.08\ \text{m/s}
\]

\textbf{Design dive speed} (FAR 23 No.23.335):

\[
V_D = 1.40 \times V_c = 1.40 \times 6.39 = 8.95\ \text{m/s}
\]

\textbf{Key V-n envelope points:}

\begin{center}
\begin{tabular}{|l|l|l|l|}
\hline
Point & Speed (m/s) & Load factor n & Condition \\
\hline
A & 8.08 & +2.5 & Positive manoeuvre limit \\
\hline
D & 8.95 & +2.5 & Design dive speed (positive) \\
\hline
G & 8.08 & -1.0 & Negative manoeuvre limit \\
\hline
Gust (cruise, unalleviated) & 6.39 & +8.31 & ${\Delta}$n = 7.31 at V\_c \\
\hline
Gust (cruise, alleviated) & 6.39 & +3.43 & K\_g ${\times}$ ${\Delta}$n at V\_c \\
\hline
1 g cruise & 6.39 & +1.0 & Nominal level flight \\
\hline
\end{tabular}
\end{center}

\textbf{Design envelope summary:}
\begin{itemize}[noitemsep,topsep=2pt]
\item Positive limit load factor: \textbf{n = +2.5} (manoeuvre governs after alleviation)
\item Negative limit load factor: \textbf{n = -1.0} (conservative; solar UAVs rarely inverted)
\item Design dive speed: \textbf{V\_D = 8.95 m/s}
\end{itemize}

\begin{figure}[H]
    \centering
    \includegraphics[width=0.905\linewidth]{Fig2_V-n_Flight_Envelope_Diagram.png} 
    \caption{V-n flight envelope diagram. m = 0.5kg, w/s = 20N/m\(^2\), K\(_{\text{g}}\) = 0.332.}
    \label{fig:fig2}
\end{figure}


\subsection{Aileron Roll Manoeuvre Load}


For the \textbf{asymmetric roll manoeuvre} (one aileron deflected ${\approx}$ 20$^{\circ}$), one wing generates increased lift and the other reduced lift. This constitutes an additional load case.

Aileron effectiveness assumption: 10 \% incremental lift per wing relative to steady-flight value.

\[
\Delta L_{\text{aileron}} = \pm 0.10 \times L_{1/2,\text{limit}} = \pm 0.10 \times 6.13 = \pm 0.613\ \text{N}
\]

\begin{center}
\begin{tabular}{|l|l|l|}
\hline
Wing & Lift at limit, roll manoeuvre (N) &  \\
\hline
Up-going (loaded) & 6.13 + 0.61 = \textbf{6.74 N} & +10 \% increment \\
\hline
Down-going (unloaded) & 6.13 - 0.61 = \textbf{5.52 N} & -10 \% increment \\
\hline
\end{tabular}
\end{center}

The up-going wing governs: $L_{1/2,\text{limit,roll}} = 6.74\ \text{N}$, which is \textbf{+10 \% above the symmetric limit load} and remains within the manoeuvre envelope. The rear spar (aileron hinge line) sees an additional local concentrated hinge-load increment.

\textbf{Rear spar incremental root bending moment from aileron:}

\[
\Delta M_{\text{RS,root}} = \Delta L_{\text{aileron}} \times \bar{y}_{\text{aileron}}
\]

Aileron assumed to span the outer 30 \% of semi-span: $\bar{y}_{\text{aileron}} = 0.85 \times (b/2) = 0.85 \times 0.490 = 0.417\ \text{m}$

\[
\Delta M_{\text{RS,root}} = 0.613 \times 0.417 = 0.256\ \text{N}\cdot\text{m}
\]

From Section 8.2, the rear spar root bending moment at symmetric limit is M\_RS,limit = 0.191 N${\cdot}$m (10 \% of total, from No.9.1). With aileron increment:

\[
M_{\text{RS,root,roll}} = 0.191 + 0.256 = 0.447\ \text{N}\cdot\text{m}
\]

Checking the rear spar ($\varnothing$3 mm rod, I = 3.976 ${\times}$ 10$^{-}$$^{1}$$^{2}$ m$^{4}$, ${\sigma}$\_allow = 206.7 MPa):

\[
\sigma_{\text{RS,roll}} = \frac{M_{\text{RS,root,roll}} \times r}{I} = \frac{0.447 \times 1.5 \times 10^{-3}}{3.976 \times 10^{-12}} = \frac{6.705 \times 10^{-4}}{3.976 \times 10^{-12}} = 168.6\ \text{MPa}
\]

\[
\text{MS}_{\text{RS,roll}} = \frac{206.7}{168.6} - 1 = +0.23 \quad \checkmark
\]

The rear spar survives the aileron roll load case with a positive margin of +0.23. $\checkmark $


\subsection{Aerodynamic Loads at Operating Altitude (ISA 1000 m)}


Solar UAVs frequently cruise above the ground boundary layer to benefit from steadier winds and improved solar irradiance. A representative operating altitude of \textbf{1000 m AMSL} is analysed here using the International Standard Atmosphere (ISA) model.


\subsubsection{ISA atmospheric properties at 1000 m}


\[
T_{1000} = 288.15 - 6.5 \times 1.0 = 281.65\ \text{K}
\]

\[
p_{1000} = 101{,}325 \times \left(1 - \frac{0.0065 \times 1000}{288.15}\right)^{g/(R_{\text{air}} \cdot \Lambda)} = 101{,}325 \times (0.9774)^{5.256} = 89{,}876\ \text{Pa}
\]

where the exponent $g/(R_{\text{air}} \cdot \Lambda) = 9.807/(287.05 \times 0.0065) = 5.256$ ($\Lambda = 0.0065\ \text{K/m}$ is the ISA tropospheric temperature lapse rate; ICAO Doc 7488).

\[
\rho_{1000} = \frac{p_{1000}}{R\,T_{1000}} = \frac{89{,}876}{287.05 \times 281.65} = \frac{89{,}876}{80{,}839} = 1.112\ \text{kg/m}^3
\]

\[
\mu_{1000} \approx 1.758 \times 10^{-5}\ \text{Pa}\cdot\text{s} \quad (\text{from Sutherland's law at } 281.65\ \text{K})
\]

\begin{center}
\begin{tabular}{|l|l|l|l|}
\hline
Quantity & Sea Level (ISA) & 1000 m (ISA) & Change \\
\hline
Temperature T & 288.15 K & 281.65 K & -2.3 \% \\
\hline
Pressure p & 101,325 Pa & 89,876 Pa & -11.3 \% \\
\hline
Air density ${\rho}$ & 1.225 kg/m$^{3}$ & 1.112 kg/m$^{3}$ & -9.2 \% \\
\hline
Dynamic viscosity ${\mu}$ & 1.789 ${\times}$ 10$^{-}$$^{5}$ Pa${\cdot}$s & 1.758 ${\times}$ 10$^{-}$$^{5}$ Pa${\cdot}$s & -1.7 \% \\
\hline
\end{tabular}
\end{center}


\subsubsection{Cruise speed at 1000 m}


At altitude the wing must still produce the same lift to support the aircraft weight. The cruise speed increases to compensate for lower density:

\[
v_{1000} = \sqrt{\frac{2L}{\rho_{1000}\,C_L\,S}} = \sqrt{\frac{2 \times 4.905}{1.112 \times 0.80 \times 0.245}} = \sqrt{\frac{9.810}{0.2180}} = \sqrt{45.00} = 6.71\ \text{m/s}
\]

Compared with sea-level cruise speed of 6.39 m/s, this is a \textbf{+5 \% increase} - consistent with the density ratio: $v_{1000}/v_{SL} = \sqrt{\rho_{SL}/\rho_{1000}} = \sqrt{1.225/1.112} = 1.050$.


\subsubsection{Dynamic pressure and Reynolds number at 1000 m}


\[
q_{\infty,1000} = \frac{1}{2}\,\rho_{1000}\,v_{1000}^2 = \frac{1}{2} \times 1.112 \times (6.71)^2 = 0.556 \times 45.02 = 25.03\ \text{N/m}^2
\]

Dynamic pressure at 1000 m (25.03 N/m$^{2}$) is virtually identical to the sea-level value (25.01 N/m$^{2}$). This is expected - for level flight the required lift equals weight, so $q_\infty\,C_L\,S = W$ regardless of altitude, meaning $q_\infty$ is fixed by the aircraft weight and wing parameters, not by altitude directly.

Reynolds number at 1000 m:

\[
Re_{1000} = \frac{\rho_{1000}\,v_{1000}\,c}{\mu_{1000}} = \frac{1.112 \times 6.71 \times 0.250}{1.758 \times 10^{-5}} = \frac{1.866}{1.758 \times 10^{-5}} = 1.06 \times 10^5
\]

The Reynolds number decreases slightly relative to sea level (1.09 ${\times}$ 10$^{5}$), remaining within the valid operating envelope of the WE3.55-9.3 airfoil (80000 - 300000). Transition behaviour may shift marginally, but no change in structural sizing is required.


\subsubsection{Aerodynamic loads at 1000 m - summary}


Because $q_\infty$ is essentially constant with altitude at constant $C_L$ flight, the \textbf{aerodynamic forces and structural loads are unchanged} from the sea-level analysis.

\begin{center}
\begin{tabular}{|l|l|l|l|}
\hline
Load & Sea Level & 1000 m &  \\
\hline
Cruise $q_\infty$ & 25.01 N/m$^{2}$ & 25.03 N/m$^{2}$ & Negligible change \\
\hline
Cruise speed & 6.39 m/s & 6.71 m/s & +5 \% \\
\hline
Total lift L & 4.905 N & 4.905 N & = W (level flight) \\
\hline
Limit lift $L_{\text{limit}}$ & 12.26 N & 12.26 N & n = 2.5 unchanged \\
\hline
Ultimate lift $L_{\text{ult}}$ & 18.39 N & 18.39 N & FOS = 1.5 unchanged \\
\hline
Half-wing shear $V_{\text{root,ult}}$ & 9.20 N & 9.20 N & Unchanged \\
\hline
Root bending moment $M_{\text{root,ult}}$ & 1.913 N${\cdot}$m & 1.913 N${\cdot}$m & Unchanged \\
\hline
Root torque $T_{\text{root}}$ & 0.105 N${\cdot}$m & 0.105 N${\cdot}$m & Unchanged \\
\hline
\end{tabular}
\end{center}

\textbf{Gust load at 1000 m:}

The gust load increment is sensitive to density. Using the CS-23 formula with $\rho_{1000} = 1.112\ \text{kg/m}^3$:

\[
\Delta n_{\text{gust},1000} = \frac{\rho_{1000}\,U_e\,v_{1000}\,a}{2(W/S)} = \frac{1.112 \times 9.1 \times 6.71 \times 4.16}{2 \times 20.0} = \frac{284.5}{40.0} = 7.11
\]

The gust load increment at 1000 m is \textbf{${\Delta}$n\_gust = 7.11}, compared with 7.31 at sea level - a \textbf{~3 \% reduction} due to slightly lower air density at altitude. This increment is an \textit{additional} load above the 1-g cruise condition; i.e., a gust encounter at cruise would produce a total load factor of 1 + 7.11 = \textbf{8.11 g}, well beyond the structural limit of n = 2.5. The governing structural design load therefore remains the manoeuvre limit n = 2.5, and operationally, flight must be restricted to calm conditions (gusts < 5 m/s) to keep the effective load factor within the structural envelope.

\textbf{Design conclusion:} All structural sizing in sections 5-12 remains valid up to 1000 m operating altitude. No altitude-driven rescaling of spar, rib, or skin dimensions is required.

\newpage \section{Distributed Lift and Shear Force Diagrams}

The wing coordinate y runs from root (y = 0) to tip (y = b/2 = 0.490 m).

\subsection{Uniform load (rectangular wing, conservative)}

Lift per unit span (uniform assumption):
\[
w_0 = \frac{L_{1/2,\text{ult}}}{b/2} = \frac{9.20}{0.490} = 18.78\ \text{N/m}
\]

\subsection{Elliptical load distribution (Prandtl lifting-line correction)}

For a rectangular planform at low AR, the true distribution lies between uniform and elliptical. The elliptical distribution is:
\[
w(y) = w_{\text{root}} \sqrt{1 - \left(\frac{y}{b/2}\right)^2}, \quad y \in [0,\ b/2]
\]

where the root value is found from the total half-span lift:

\[
L_{1/2} = \int_0^{b/2} w(y)\, dy = w_{\text{root}} \int_0^{b/2} \sqrt{1 - \left(\frac{y}{b/2}\right)^2}\, dy = w_{\text{root}} \times \frac{\pi (b/2)}{4}
\]

\[
w_{\text{root}} = \frac{4\, L_{1/2,\text{ult}}}{\pi\, (b/2)} = \frac{4 \times 9.20}{\pi \times 0.490} = \frac{36.80}{1.539} = 23.91\ \text{N/m}
\]

Spot values (ultimate):

\begin{center}
\begin{tabular}{|l|l|l|}
\hline
y (m) & y/(b/2) & w(y) (N/m) \\
\hline
0.000 (root) & 0.000 & 23.91 \\
\hline
0.098 & 0.200 & 23.43 \\
\hline
0.196 & 0.400 & 21.90 \\
\hline
0.294 & 0.600 & 19.12 \\
\hline
0.392 & 0.800 & 14.34 \\
\hline
0.441 & 0.900 & 10.41 \\
\hline
0.490 (tip) & 1.000 & 0.00 \\
\hline
\end{tabular}
\end{center}

\begin{figure}[H]
    \centering
    \includegraphics[width=0.905\linewidth]{Fig3_Lift_Distribution_(Elliptical_vs_Uniform).png}
    \caption{Distributed lift load distribution (Ultimate loads). L\(_\text{1/2,ult}\) = 9.20N, n = 2.5 and FOS = 1.5.}
    \label{fig:fig3}
\end{figure}


\subsection{Shear force diagram V(y) - ultimate loads}

The shear force at station y (looking from tip toward root - cantilever convention, positive up):

\[
V(y) = \int_y^{b/2} w(\xi)\, d\xi
\]

For the elliptical distribution:

\[
V(y) = w_{\text{root}} \int_y^{b/2} \sqrt{1 - \left(\frac{\xi}{b/2}\right)^2}\, d\xi
\]

Let ${\eta}$ = $\xi$/(b/2), d$\xi$ = (b/2) d${\eta}$:

\[
V(y) = w_{\text{root}} \frac{b/2}{2} \left[\eta\sqrt{1-\eta^2} + \arcsin\eta\right]_{\eta_y}^{1}
\]

where ${\eta}$\_y = y/(b/2) and the upper limit evaluates to ${\pi}$/2.

At root (y = 0, ${\eta}$\_y = 0):

\[
V(0) = w_{\text{root}} \frac{b/2}{2} \left[\frac{\pi}{2} - 0\right] = 23.91 \times \frac{0.490}{2} \times \frac{\pi}{2}
\]

\[
= 23.91 \times 0.245 \times 1.5708 = 9.20\ \text{N} \quad (\checkmark = L_{1/2,\text{ult}})
\]

Spot values of V(y) (ultimate):

\begin{center}
\begin{tabular}{|l|l|}
\hline
y (m) & V(y) (N) \\
\hline
0.000 (root) & 9.20 \\
\hline
0.098 & 7.65 \\
\hline
0.196 & 5.87 \\
\hline
0.294 & 3.91 \\
\hline
0.392 & 1.88 \\
\hline
0.441 & 0.97 \\
\hline
0.490 (tip) & 0.00 \\
\hline
\end{tabular}
\end{center}

\begin{figure}[H]
    \centering
    \includegraphics[width=0.905\linewidth]{Fig4_Shear_Force_Diagram_V(y).png}
    \caption{Shear force Diagram V(y) (Ultimate Loads).}
    \label{fig:}
\end{figure}

\subsection{Inertia Relief}

The distributed weight of the wing structure acts downward during flight, partially opposing the upward aerodynamic lift and reducing the net structural loads at the root. This is called \textbf{inertia relief}.

\textbf{Wing structural mass per unit span:}

\[
\bar{m} = \frac{m_{\text{wing,half}}}{b/2}
\]

From Section 13.4, total wing structure (both halves) ${\approx}$ 203 g; per half-span: $m_{\text{wing,half}} \approx 101.5\ \text{g}$.

\[
\bar{m} = \frac{0.1015}{0.490} = 0.207\ \text{kg/m}
\]

\textbf{Inertia load per unit span (uniform, downward at load factor n):}

\[
w_{\text{inertia}}(y) = \bar{m} \times g \times n = 0.207 \times 9.81 \times 2.5 = 5.08\ \text{N/m}
\]

\textbf{Net aerodynamic load (elliptical lift minus inertia relief):}

\[
q_{\text{net}}(y) = w(y) - w_{\text{inertia}} = 23.91\sqrt{1-(y/0.490)^2} - 5.08 \quad [\text{N/m at ultimate loads}]
\]

\begin{center}
\begin{tabular}{|l|l|l|l|}
\hline
y (m) & w(y) lift (N/m) & w\_inertia (N/m) & q\_net(y) (N/m) \\
\hline
0.000 (root) & 23.91 & 5.08 & 18.83 \\
\hline
0.196 & 21.90 & 5.08 & 16.82 \\
\hline
0.392 & 14.34 & 5.08 & 9.26 \\
\hline
0.490 (tip) & 0.00 & 5.08 & -5.08 (self-weight only) \\
\hline
\end{tabular}
\end{center}

\textbf{Inertia-relieved root shear and bending moment:}

\[
V_{\text{net}}(0) = V_{\text{lift}}(0) - w_{\text{inertia}} \times (b/2) = 9.20 - 5.08 \times 0.490 = 9.20 - 2.49 = 6.71\ \text{N}
\]

\[
M_{\text{net,root}} = M_{\text{lift,root}} - \frac{w_{\text{inertia}} \times (b/2)^2}{2} = 2.025 - \frac{5.08 \times 0.240}{2} = 2.025 - 0.610 = 1.415\ \text{N}\cdot\text{m}
\]

The inertia relief reduces the root bending moment by \textbf{~30 \%} relative to the aerodynamic-only value. For structural sizing the \textbf{conservative unreduced value (M\_root,ult = 2.025 N${\cdot}$m)} is used (consistent with No.6 and No.7 design calculations), making inertia relief an additional safety margin.

\begin{figure}[H]
    \centering
    \includegraphics[width=0.85\linewidth]{Fig5_Inertia-Relief_Net_Load_Distribution.png}
    \caption{Inertia relieved net distribution load (Ultimate).}
    \label{fig:fig5}
\end{figure}

\newpage \section{Bending Moment Distribution}

\subsection{Bending moment M(y) - ultimate loads}

The bending moment at station y (root most critical, M is positive for upward bending = tension on lower surface):
\[
M(y) = \int_y^{b/2} w(\xi)(\xi - y)\, d\xi
\]

For the elliptical distribution, the moment at the root (y = 0) is:
\[
M_{\text{root}} = \int_0^{b/2} w(\xi)\, \xi\, d\xi = w_{\text{root}} \int_0^{b/2} \xi\sqrt{1 - \left(\frac{\xi}{b/2}\right)^2}\, d\xi
\]

Let ${\eta}$ = $\xi$/(b/2):
\begin{align*}
M_{\text{root}} &= w_{\text{root}} (b/2)^2 \int_0^{1} \eta\sqrt{1-\eta^2}\, d\eta = w_{\text{root}} (b/2)^2 \times \frac{1}{3}\\
M_{\text{root}} &= 23.91 \times (0.490)^2 \times \frac{1}{3} = 23.91 \times 0.2401 \times 0.3333 = 1.913\ \text{N}\cdot\text{m}
\end{align*}

Cross-check with centroid of load:

\[
\bar{y} = \frac{M_{\text{root}}}{V_{\text{root}}} = \frac{1.913}{9.20} = 0.208\ \text{m} \quad \text{($\approx$  42 \% of half-span - typical for elliptical distribution)}
\]

For uniform distribution (conservative check):
\[
M_{\text{root,uniform}} = \frac{w_0\, (b/2)^2}{2} = \frac{18.78 \times 0.490^2}{2} = \frac{18.78 \times 0.2401}{2} = 2.253\ \text{N}\cdot\text{m}
\]

\textbf{Design bending moment at root (ultimate):}
\[
\boxed{M_{\text{root,ult}} = 1.913\ \text{N}\cdot\text{m} \quad \text{(elliptical)}}
\]

Use $M_{\text{design}} = 2.25\ \text{N}\cdot\text{m}$ (uniform, 18\% higher) for a conservative design.

Bending moment at intermediate stations (elliptical):
\[
M(y) = w_{\text{root}} \frac{(b/2)^2}{3}\left(1 - \frac{3}{2}\left(\frac{y}{b/2}\right)^2 + \frac{1}{2}\left(\frac{y}{b/2}\right)^3\right)\ \text{(approx.)}
\]

\begin{center}
\begin{tabular}{|l|l|l|}
\hline
y (m) & ${\eta}$ = y/(b/2) & M(y) (N${\cdot}$m) \\
\hline
0.000 (root) & 0.000 & 1.913 \\
\hline
0.098 & 0.200 & 1.526 \\
\hline
0.196 & 0.400 & 1.043 \\
\hline
0.294 & 0.600 & 0.573 \\
\hline
0.392 & 0.800 & 0.177 \\
\hline
0.441 & 0.900 & 0.057 \\
\hline
0.490 (tip) & 1.000 & 0.000 \\
\hline
\end{tabular}
\end{center}

\begin{figure}[H]
    \centering
    \includegraphics[width=0.85\linewidth]{Fig6_Bending_Moment_Distribution_M(y).png}
    \caption{Bending moment distribution M(y) (Ultimate loads).}
    \label{fig:fig6}
\end{figure}

\subsection{Section modulus requirement}

The main spar carries approximately 90\% of the bending moment (as established for high-AR sailplane-style wings; Noth IDR):

\[
M_{spar} = 0.90 \times M_{\text{design}} = 0.90 \times 2.25 = 2.025\ \text{N}\cdot\text{m}
\]

\newpage \section{Main Spar (Front Spar) Sizing}

The main spar is a \textbf{solid circular aluminium rod} (Al 6061-T6).

\subsection{Material properties - Al 6061-T6}

\begin{center}
\begin{tabular}{|l|l|l|}
\hline
Property & Symbol & Value \\
\hline
Young's modulus & E & 69 GPa = 69 ${\times}$ 10$^{9}$ Pa \\
\hline
Shear modulus & G & 26 GPa = 26 ${\times}$ 10$^{9}$ Pa \\
\hline
Yield strength & ${\sigma}$\_y & 276 MPa \\
\hline
Ultimate tensile strength & ${\sigma}$\_ult & 310 MPa \\
\hline
Density & ${\rho}$\_Al & 2700 kg/m$^{3}$ \\
\hline
Allowable bending stress (= ${\sigma}$\_ult / FOS) & ${\sigma}$\_allow & 310 / 1.5 = \textbf{206.7 MPa} \\
\hline
Allowable shear stress (= 0.577 ${\times}$ ${\sigma}$\_y) & ${\tau}$\_allow & 0.577 ${\times}$ 276 = 159.2 MPa \\
\hline
\end{tabular}
\end{center}

\subsection{Minimum spar radius from bending}

Solid circular rod, bending about a diameter.  
Section modulus: $Z = I/y = (\pi r^4/4)/r = \pi r^3/4$

\begin{align*}
\sigma &= \frac{M}{Z} = \frac{4M}{\pi r^3} \leq \sigma_{\text{allow}}\\
r^3 &\geq \frac{4 M_{spar}}{\pi \sigma_{\text{allow}}} = \frac{4 \times 2.025}{\pi \times 206.7 \times 10^6} = \frac{8.10}{649.1 \times 10^6} = 1.248 \times 10^{-8}\ \text{m}^3 \\
r &\geq (1.248 \times 10^{-8})^{1/3} = 2.32 \times 10^{-3}\ \text{m} = 2.32\ \text{mm} \\
\end{align*}

Minimum diameter: $d_{\min} = 2 \times 2.32 = 4.64\ \text{mm}$

We select a standard rod: $d_{FS} = 5\ \text{mm}$, $r_{FS} = 2.5\ \text{mm}$

Actual peak bending stress:
\[
\sigma_{FS} = \frac{4 \times 2.025}{\pi \times (2.5 \times 10^{-3})^3} = \frac{8.10}{\pi \times 1.5625 \times 10^{-8}} = \frac{8.10}{4.909 \times 10^{-8}} = 165.0\ \text{MPa}
\]

Margin of safety: $MS = \sigma_{\text{allow}}/\sigma_{FS} - 1 = 206.7/165.0 - 1 = +0.25$ $\checkmark $

\subsection{Shear check on main spar}

Maximum transverse shear at root:

\[
V_{\text{root,ult}} = 9.20\ \text{N}
\]

For solid circular section, maximum shear stress:

\begin{align*}
\tau_{\max} &= \frac{4V}{3A} \\
&= \frac{4 \times 9.20}{3 \times \pi (2.5 \times 10^{-3})^2} \\
&= \frac{36.80}{3 \times \pi \times 6.25 \times 10^{-6}} \\
&= \frac{36.80}{58.905 \times 10^{-6}} \\
&= 0.625\ \text{MPa} \\
\\
\tau_{\max} &\ll \tau_{\text{allow}} = 159.2\ \text{MPa} 
\quad (\text{shear is not critical})
\end{align*}

\subsection{Spar deflection at tip}

Beam equation for cantilever with elliptical distributed load. The tip deflection for a uniform load (conservative):

\[
\delta_{\text{tip}} = \frac{w_0 (b/2)^4}{8 EI}
\]

Second moment of area:

\begin{align*}
I_{FS} &= \frac{\pi r^4}{4} \\
&= \frac{\pi (2.5 \times 10^{-3})^4}{4} \\
&= \frac{\pi \times 3.90625 \times 10^{-11}}{4} \\
&= 3.068 \times 10^{-11}\ \text{m}^4 \\
\\
\delta_{\text{tip}} &= \frac{18.78 \times (0.490)^4}{8 \times 69 \times 10^9 \times 3.068 \times 10^{-11}} \\
&= \frac{18.78 \times 0.05765}{8 \times 2.117 \times 10^{-3}} \\
&= \frac{1.0826}{1.6936 \times 10^{-2}} \\
&= 0.0639\ \text{m} = 63.9\ \text{mm}
\end{align*}

This deflection (${\approx}$ 6.5 \% of semi-span) is high for a 5 mm rod. It is acceptable for a flexible solar UAV wing but confirms the need for careful aeroelastic assessment. The deflection under limit loads (n = 2.5, before FOS):

$$\delta\_{\text{tip,limit}} = \frac{63.9}{FOS} = \frac{63.9}{1.5} = 42.6\ \text{mm} \approx 4.4\%\ \text{semi-span}$$

This is within the typical 5\% limit for flexible UAV wings.

\subsection{Spar mass}
\begin{align*}
m_{FS} &= \rho_{Al} \, A_{FS} \, \frac{b}{2} \\
&= 2700 \times \pi (2.5 \times 10^{-3})^2 \times 0.490 \\
&= 2700 \times 1.9635 \times 10^{-5} \times 0.490 \\
&= 2700 \times 9.621 \times 10^{-6} \\
&= 0.02598\ \text{kg} = 25.98\ \text{g} \quad \text{(per half-spar)} \\
\\
m_{FS,\text{total}} &= 2 \times 25.98 = 51.96\ \text{g}
\end{align*}

\begin{figure}[H]
    \centering
    \includegraphics[width=0.905\linewidth]{Fig7_Spar_Deflection_Shape_Along_Span.png}
    \caption{Font spar deflection shape along semi-span.}
    \label{fig:fig7}
\end{figure}

\newpage \section{Rear Spar Sizing}

The rear spar carries approximately 10\% of the bending moment and handles hinge loads from the ailerons plus torsion.
\[
M_{RS} = 0.10 \times 2.25 = 0.225\ \text{N}\cdot\text{m}
\]

\textbf{Minimum radius:}
\[
r_{RS} \geq \left(\frac{4 \times 0.225}{\pi \times 206.7 \times 10^6}\right)^{1/3} = \left(\frac{0.900}{649.1 \times 10^6}\right)^{1/3} = (1.386 \times 10^{-9})^{1/3} = 1.115 \times 10^{-3}\ \text{m} = 1.12\ \text{mm}
\]

Select **standard rod: $d_{RS} = 3\ \text{mm}$, $r_{RS} = 1.5\ \text{mm}$** (nearest standard size above minimum).

Actual stress:
\[
\sigma_{RS} = \frac{4 \times 0.225}{\pi (1.5 \times 10^{-3})^3} = \frac{0.900}{\pi \times 3.375 \times 10^{-9}} = \frac{0.900}{1.060 \times 10^{-8}} = 84.9\ \text{MPa}
\]

Margin: $MS = 206.7/84.9 - 1 = +1.43$ $\checkmark $ (well within limits, but 3 mm is the minimum practical rod size)

\textbf{Rear spar mass:}
\[
m_{RS,\text{total}} = 2 \times 2700 \times \pi (1.5 \times 10^{-3})^2 \times 0.490 = 2 \times 2700 \times 7.069 \times 10^{-6} \times 0.490 = 18.74\ \text{g}
\]

\newpage \section{Rib Design and Spacing}

\subsection{Number of ribs and spacing}

Ribs prevent local skin buckling and maintain the airfoil shape (especially critical for solar-cell adhesion on the flat upper surface).

Total ribs:
\[
N_{\text{ribs}} = 40 \quad \text{(per design brief)}
\]

Rib pitch (uniform):
\[
\Delta y = \frac{b}{N_{\text{ribs}} - 1} = \frac{980}{39} \approx 25.1\ \text{mm}
\]

Ribs are placed at: $y_i = (i-1) \times 25.1\ \text{mm}$, for $i = 1, 2, \ldots, 40$.

\subsection{Rib load}

Each rib distributes the aerodynamic load from the skin to the spars. The rib carries the distributed load over one rib bay:
\[
F_{\text{rib}}(y_i) = w(y_i) \times \Delta y
\]

At root (maximum):
\[
F_{\text{rib,root}} = w(0) \times \Delta y = 23.91\ \text{N/m} \times 0.0251\ \text{m} = 0.600\ \text{N (per rib at root)}
\]

\subsection{Rib as a simply-supported beam}

Each rib spans between the leading edge and the rear spar (chord-wise beam), loaded by the distributed aerodynamic pressure over its tributary area. The rib is modelled as a simply supported beam between the two spars (main structural load path).

Effective span of rib: $L_{\text{rib}} = x_{RS} - x_{FS} = 162.5 - 62.5 = 100\ \text{mm} = 0.100\ \text{m}$

Distributed load on rib (pressure over tributary area):
\[
q_{\text{rib}} = \frac{F_{\text{rib}}}{L_{\text{rib}}} = \frac{0.600}{0.100} = 6.0\ \text{N/m}
\]

Maximum bending moment in rib (SS beam, UDL):
\[
M_{\text{rib,max}} = \frac{q_{\text{rib}}\, L_{\text{rib}}^2}{8} = \frac{6.0 \times (0.100)^2}{8} = \frac{6.0 \times 0.01}{8} = 7.5 \times 10^{-3}\ \text{N}\cdot\text{m}
\]

\subsection{Rib cross-section sizing (balsa wood)}

\textbf{Balsa wood properties (medium density):}

\begin{center}
\begin{tabular}{|l|l|l|}
\hline
Property & Symbol & Value \\
\hline
Young's modulus (along grain) & E\_b & 3.4 GPa \\
\hline
Bending strength (MOR) & ${\sigma}$\_b & 14.7 MPa \\
\hline
Density & ${\rho}$\_b & 160 kg/m$^{3}$ \\
\hline
Allowable bending stress (= ${\sigma}$\_b / FOS) & ${\sigma}$\_allow,b & 14.7 / 1.5 = \textbf{9.8 MPa} \\
\hline
\end{tabular}
\end{center}

Rib cross-section: rectangular \textbf{width b\_r ${\times}$ depth h\_r} (flat sheet rib).

The rib depth is constrained by the airfoil profile. At mid-span of the rib (between FS and RS, at 45\% chord ${\approx}$ 112 mm from LE), the airfoil height is approximately:
\[
h_r \approx 0.90 \times t_{\max} = 0.90 \times 23.25 = 20.9\ \text{mm}
\]

\textbf{Section modulus of rectangular section}: $Z_r = b_r h_r^2 / 6$
\[
b_r \geq \frac{6 M_{\text{rib,max}}}{\sigma_{\text{allow,b}}\, h_r^2} = \frac{6 \times 7.5 \times 10^{-3}}{9.8 \times 10^6 \times (20.9 \times 10^{-3})^2}
\]
\[
= \frac{0.045}{9.8 \times 10^6 \times 4.368 \times 10^{-4}} = \frac{0.045}{4280.6} = 1.051 \times 10^{-5}\ \text{m} = 0.011\ \text{mm}
\]

This is essentially zero - the ribs are sized by \textbf{minimum practical thickness}, not by stress. Use:
\[
b_r = t_{\text{rib}} = 2.0\ \text{mm} \quad \text{(standard balsa sheet thickness)}
\]

\textbf{Check}: $\sigma_r = 6 M / (b_r h_r^2) = 6 \times 7.5 \times 10^{-3} / (0.002 \times (0.0209)^2) = 0.045 / (8.737 \times 10^{-7}) = 51,510\ \text{Pa} \approx 0.05\ \text{MPa}$

$\sigma_r \ll \sigma_{\text{allow,b}} = 9.8\ \text{MPa}$ $\checkmark $ - ribs are very lightly loaded.

\subsection{Skin-buckling check (rib spacing)}

The critical buckling stress of a simply-supported flat plate (skin panel between ribs) under compression is (Roark's Formulas, flat plate buckling):
\[
\sigma_{\text{cr}} = k \frac{\pi^2 E_s}{12(1-\nu^2)} \left(\frac{t_s}{a}\right)^2
\]

where:
\begin{itemize}[noitemsep,topsep=2pt]
\item $k$ = 4.0 (SS on all edges, compression in chord direction)
\item $E_s$ = 69 GPa (aluminium)
\item $\nu$ = 0.33 (aluminium)
\item $t_s$ = skin thickness (to be determined in No.11)
\item $a$ = rib spacing = 25.1 mm
\end{itemize}

Setting $\sigma_{\text{cr}} = \sigma_{\text{allow}}$ and solving for $t_s$ will be done in No.11.

\subsection{Rib mass}

Each rib is a thin balsa sheet of approximate plan area A\_rib:
\[
A_{\text{rib}} \approx c \times t_{\max} \times 0.60 = 0.250 \times 0.02325 \times 0.60 \approx 3.49 \times 10^{-3}\ \text{m}^2 \quad (\text{approx. airfoil cross-section area})
\]

More accurately, rib area ${\approx}$ airfoil chord ${\times}$ average height ${\times}$ 50\%:
\[
A_{\text{rib}} \approx 0.250 \times 0.020 / 2 = 2.5 \times 10^{-3}\ \text{m}^2
\]

Mass of one rib:
\[
m_{\text{one rib}} = \rho_b \times t_{\text{rib}} \times A_{\text{rib}} = 160 \times 0.002 \times 2.5 \times 10^{-3} = 8.0 \times 10^{-4}\ \text{kg} = 0.80\ \text{g}
\]

Total rib mass:
\[
m_{\text{ribs}} = 40 \times 0.80 = 32.0\ \text{g}
\]

\newpage \section{Alternative Rib Design: Acrylic Sheet Ribs}

This section analyses an alternative rib material --- \textbf{cast acrylic (PMMA)} --- and determines the required rib thickness and its effect on the overall structural design. All spar sizing remains unchanged; only the rib material changes.

\subsection{Acrylic (PMMA) Material Properties}

\begin{center}
\begin{tabular}{|l|l|l|}
\hline
Property & Symbol & Value \\
\hline
Young's modulus & E\_ac & 3.2 GPa \\
\hline
Flexural strength (MOR) & ${\sigma}$\_f & 100 MPa \\
\hline
Tensile strength & ${\sigma}$\_t & 70 MPa \\
\hline
Density & ${\rho}$\_ac & 1180 kg/m$^{3}$ \\
\hline
Allowable bending stress ($= \sigma_f / \text{FOS}$) & ${\sigma}$\_allow,ac & 100 / 1.5 = \textbf{66.7 MPa} \\
\hline
Shear modulus & G\_ac & $\approx$ 1.2 GPa \\
\hline
Poisson's ratio & ${\nu}$ & 0.37 \\
\hline
\end{tabular}
\end{center}

Compared with balsa wood: acrylic is \textbf{7.4$\times$ heavier} ($\rho_{ac}/\rho_b = 1180/160$) but \textbf{6.8$\times$ stronger} in bending allowable ($\sigma_{\text{allow,ac}}/\sigma_{\text{allow,b}} = 66.7/9.8$). The result is a more rigid but significantly heavier rib.

\subsection{Rib Load (Same as No.9.2)}

The rib loads are independent of rib material. Using the existing 40-rib design (No.9):

\[
\Delta y = 25.1\ \text{mm}, \quad F_{\text{rib,root}} = 23.91 \times 0.0251 = 0.600\ \text{N}
\]

\[
L_{\text{rib}} = 100\ \text{mm}, \quad q_{\text{rib}} = 6.0\ \text{N/m}
\]

\[
M_{\text{rib,max}} = \frac{q_{\text{rib}}\, L_{\text{rib}}^2}{8} = \frac{6.0 \times (0.100)^2}{8} = 7.5 \times 10^{-3}\ \text{N}\cdot\text{m}
\]

Rib depth at mid-span (45\% chord, WE3.55-9.3):
\[
h_r \approx 0.90 \times t_{\max} = 0.90 \times 23.25 = 20.9\ \text{mm}
\]

\subsection{Acrylic Rib Thickness Calculation}

\textbf{Section modulus:} $Z_r = t_{\text{rib}} \times h_r^2 / 6$

Setting $\sigma = M/Z \leq \sigma_{\text{allow,ac}}$:

\[
t_{\text{rib}} \geq \frac{6 M_{\text{rib,max}}}{\sigma_{\text{allow,ac}}\, h_r^2} = \frac{6 \times 7.5 \times 10^{-3}}{66.7 \times 10^6 \times (20.9 \times 10^{-3})^2}
\]

\[
= \frac{0.045}{66.7 \times 10^6 \times 4.368 \times 10^{-4}} = \frac{0.045}{29{,}134} = 1.545 \times 10^{-6}\ \text{m} = 0.0015\ \text{mm}
\]

The stress-derived minimum is negligible. \textbf{Practical minimum for laser-cut acrylic: 2.0 mm.}

\[
\boxed{t_{\text{rib,acrylic}} = 2.0\ \text{mm}}
\]

\textbf{Stress check} (2 mm acrylic):
\[
\sigma_{r,ac} = \frac{6 \times 7.5 \times 10^{-3}}{0.002 \times (0.0209)^2} = \frac{0.045}{8.737 \times 10^{-7}} = 51{,}510\ \text{Pa} = 0.052\ \text{MPa}
\]

\[
\text{MS}_{ac} = \frac{66.7}{0.052} - 1 = 1282 \gg 1 \quad \checkmark
\]

\subsection{Spar Thickness with Acrylic Ribs}

Spar sizing is governed by wing bending loads, which are independent of rib material. The same aluminium alloy (Al 6061-T6) spar diameters apply:

\textbf{Front spar minimum radius from bending:}
\[
r^3 \geq \frac{4 M_{spar}}{\pi \sigma_{\text{allow}}} = \frac{4 \times 2.025}{\pi \times 206.7 \times 10^6} = 1.248 \times 10^{-8}\ \text{m}^3 \implies r_{\min} = 2.32\ \text{mm}
\]
\[
\boxed{d_{FS} = 5\ \text{mm}} \quad \sigma_{FS} = 165.0\ \text{MPa}, \quad \text{MS} = +0.25\ \checkmark
\]

\textbf{Rear spar minimum radius:}
\[
r_{RS,\min} = \left(\frac{4 \times 0.225}{\pi \times 206.7 \times 10^6}\right)^{1/3} = 1.115 \times 10^{-3}\ \text{m} = 1.12\ \text{mm}
\]
\[
\boxed{d_{RS} = 3\ \text{mm}} \quad \sigma_{RS} = 84.9\ \text{MPa}, \quad \text{MS} = +1.43\ \checkmark
\]

Spar diameters are \textbf{unchanged} from the balsa-rib design.

\subsection{Acrylic Rib Mass}

Approximate rib plan area (WE3.55-9.3, $c = 250$ mm, $t_{\max} = 23.25$ mm):
\[
A_{\text{rib}} \approx 0.250 \times 0.020 / 2 = 2.5 \times 10^{-3}\ \text{m}^2
\]

Mass of one acrylic rib (2 mm thick):
\[
m_{\text{rib,ac}} = \rho_{ac} \times t_{\text{rib,ac}} \times A_{\text{rib}} = 1180 \times 0.002 \times 2.5 \times 10^{-3} = 5.90 \times 10^{-3}\ \text{kg} = 5.90\ \text{g}
\]

Total acrylic rib mass (40 ribs):
\[
m_{\text{ribs,ac}} = 40 \times 5.90 = 236.0\ \text{g}
\]

\subsection{Mass Comparison: Balsa vs Acrylic}

\begin{center}
\begin{tabular}{|l|l|l|l|}
\hline
Component & Balsa Ribs (2 mm) & Acrylic Ribs (2 mm) & Difference \\
\hline
Front spar (both half-spars) & 52 g & 52 g & - \\
\hline
Rear spar (both half-spars) & 19 g & 19 g & - \\
\hline
Ribs (40$\times$) & 32 g & 236 g & \textbf{+204 g} \\
\hline
Leading-edge D-box skin & 50 g & 50 g & - \\
\hline
\textbf{Total structural} & \textbf{$\sim$203 g} & \textbf{$\sim$407 g} & \textbf{+204 g (+100\%)} \\
\hline
\end{tabular}
\end{center}

\subsection{Why Choose Acrylic Ribs?}

\textbf{Advantages of acrylic over balsa:}
\begin{enumerate}[noitemsep,topsep=2pt]
\item \textbf{Dimensional stability:} Acrylic does not absorb moisture; will not warp with humidity changes.
\item \textbf{Isotropy:} Balsa is strongly anisotropic (weaker across grain). Acrylic is isotropic.
\item \textbf{Precision machining:} CO$_2$ laser-cut acrylic achieves $\pm$0.1 mm tolerance on NACA profiles.
\item \textbf{Surface quality:} Smooth acrylic bonds well with structural adhesives and film covering.
\item \textbf{Impact resistance:} More tolerant of chord-wise impact (e.g., landing loads) than balsa.
\end{enumerate}

\textbf{Disadvantages:}
\begin{enumerate}[noitemsep,topsep=2pt]
\item \textbf{Mass penalty:} 236 g vs 32 g --- acrylic ribs add 204 g to structural mass.
\item \textbf{Brittleness:} Acrylic is brittle under concentrated point loads; spar notches require generous radii ($r \geq 1$ mm).
\item \textbf{Not suitable for flight-critical applications} where minimum mass is the priority.
\end{enumerate}

\textbf{Recommendation:} Use acrylic ribs for \textbf{ground test specimens}, \textbf{structural demonstration models}, or \textbf{display wings} where dimensional stability and appearance outweigh mass. For flight-critical solar UAV construction, use balsa wood (2 mm).

\newpage \section{Torsion Analysis}

\subsection{Applied aerodynamic torque}

The wing twists due to the aerodynamic pitching moment and the eccentricity of lift from the elastic axis (front spar).

\textbf{Torque per unit span} about the front spar (elastic axis):
\[
t(y) = w(y) \times e_L + m_{AC}'
\]

where:
\begin{itemize}[noitemsep,topsep=2pt]
\item $w(y)$ = distributed lift load per unit span
\item $e_L$ = 15.6 mm = 0.0156 m (lift eccentricity from FS, No.4.2)
\item $m_{AC}' = C_{M,AC} \times q \times c^2$ = pitching moment per unit span
\end{itemize}

\[
m_{AC}' = C_{M,AC} \times q \times c^2 = (-0.05) \times 25.01 \times (0.250)^2 = -0.05 \times 25.01 \times 0.0625 = -0.0782\ \text{N$\cdot$ m/m}
\]

At root (using ultimate load factor, $w_{\text{root}} = 23.91\ \text{N/m}$):
\[
t(0) = 23.91 \times 0.0156 + (-0.0782) = 0.3730 - 0.0782 = 0.2948\ \text{N$\cdot$ m/m}
\]

\textbf{Root torque (integrating elliptical distribution of lift torque):}
\begin{align*}
T_{\text{root}} 
&= \int_0^{b/2} w(y)\, e_L \, dy + m'_{AC} \times \left(\frac{b}{2}\right) \\
\\
&= e_L \times L_{1/2,\text{ult}} + m'_{AC} \times \left(\frac{b}{2}\right) \\
\\
&= 0.0156 \times 9.20 + (-0.0782) \times 0.490 \\
\\
&= 0.1435 - 0.0383 \\
\\
&= 0.1052\ \text{N}\cdot\text{m}
\end{align*}

The net torque at root is \textbf{0.105 N${\cdot}$m} (nose-up, tending to increase angle of attack on the front portion). This must be resisted by the D-box and inter-spar structure.


\subsection{D-box (closed section) torsional stiffness}


The leading-edge D-box is the primary torsion-resisting element. It is approximated as a closed thin-walled section (Bredt-Batho formula):

\[
T = 2 A_{\text{enc}} q_s \quad \Rightarrow \quad q_s = \frac{T}{2 A_{\text{enc}}}
\]

where $A_{\text{enc}}$ is the area enclosed by the D-box mid-line.

\textbf{D-box geometry (from LE to FS):}
\begin{itemize}[noitemsep,topsep=2pt]
\item Chord-wise extent: 0 to 62.5 mm
\item Approximate shape: semi-ellipse with base c$_{1}$ = 62.5 mm, height ${\approx}$ h\_FS/2 = 10 mm (semi-height of airfoil at FS)
\end{itemize}

\[
A_{\text{enc}} \approx \frac{\pi}{2} \times \frac{c_1}{2} \times \frac{h_{FS}}{2} = \frac{\pi}{2} \times 31.25 \times 10.0 = \frac{\pi}{2} \times 312.5 = 490.9\ \text{mm}^2 = 4.909 \times 10^{-4}\ \text{m}^2
\]

Shear flow in D-box skin:
\[
q_s = \frac{T_{\text{root}}}{2 A_{\text{enc}}} = \frac{0.1052}{2 \times 4.909 \times 10^{-4}} = \frac{0.1052}{9.818 \times 10^{-4}} = 107.2\ \text{N/m}
\]

Shear stress in D-box skin (skin thickness $t_s$):
\[
\tau_{D} = \frac{q_s}{t_s}
\]

For aluminium skin with $t_s = 0.3\ \text{mm}$ (refer No.11):
\[
\tau_{D} = \frac{107.2}{3 \times 10^{-4}} = 357,333\ \text{Pa} = 0.357\ \text{MPa} \ll \tau_{\text{allow}} = 159.2\ \text{MPa} \quad \checkmark 
\]

\begin{figure}[H]
    \centering
    \includegraphics[width=0.905\linewidth]{Fig9_Natural_Frequency_Mode_Shapes.png}
    \caption{Cantilever wing bending mode shapes. First two modes are shown).}
    \label{fig:fig9}
\end{figure}

\subsection{Angle of twist}

Using the Bredt-Batho formula for twist rate:
\[
\frac{d\phi}{dy} = \frac{T}{4 A_{\text{enc}}^2} \oint \frac{ds}{G t_s}
\]

For uniform skin thickness $t_s$, the perimeter of the D-box:
\[
s_D \approx \pi \sqrt{\frac{(c_1/2)^2 + (h_{FS}/2)^2}{2}} \approx \pi \sqrt{\frac{31.25^2 + 10.0^2}{2}} = \pi \sqrt{\frac{976.6 + 100}{2}} = \pi \sqrt{538.3} = \pi \times 23.2 = 72.9\ \text{mm}
\]
\[
\oint \frac{ds}{Gt_s} = \frac{s_D}{G t_s} = \frac{72.9 \times 10^{-3}}{26 \times 10^9 \times 3 \times 10^{-4}} = \frac{72.9 \times 10^{-3}}{7.8 \times 10^6} = 9.346 \times 10^{-9}\ \text{m/N}
\]

Twist rate at root (maximum torque):
\begin{align*}
\frac{d\phi}{dy}\bigg|_{\text{root}} 
&= \frac{T_{\text{root}}}{4 A_{\text{enc}}^2} \times \frac{s_D}{G t_s} \\
&= \frac{0.1052}{4 \times (4.909 \times 10^{-4})^2} \times 9.346 \times 10^{-9} \\
\\
&= \frac{0.1052}{4 \times 2.410 \times 10^{-7}} \times 9.346 \times 10^{-9} \\
\\
&= \frac{0.1052}{9.639 \times 10^{-7}} \times 9.346 \times 10^{-9} \\
\\
&= 1.091 \times 10^{5} \times 9.346 \times 10^{-9} \\
\\
&= 1.020 \times 10^{-3}\ \text{rad/m}
\end{align*}

Total twist at tip (assuming linear torque decay from root to tip):
\[
\phi_{\text{tip}} = \frac{1}{2} \times \frac{d\phi}{dy}\bigg|_{\text{root}} \times \frac{b}{2} = \frac{1}{2} \times 1.020 \times 10^{-3} \times 0.490 = 2.50 \times 10^{-4}\ \text{rad} = 0.0143^{\circ}
\]

\textbf{Aeroelastic divergence check:}

The divergence dynamic pressure is:

\[
q_D = \frac{GJ_{\text{eff}}}{e\, a\, (b/2)^2}
\]

where $GJ_{\text{eff}} = 4 A_{\text{enc}}^2 G t_s / s_D$ (Bredt closed section):
\begin{align*}
GJ_{\text{eff}} 
&= \frac{4 \times (4.909 \times 10^{-4})^2 \times 26 \times 10^9 \times 3 \times 10^{-4}}{72.9 \times 10^{-3}} \\
\\
&= \frac{4 \times 2.410 \times 10^{-7} \times 7.8 \times 10^6}{72.9 \times 10^{-3}} \\
\\
&= \frac{7.518}{72.9 \times 10^{-3}} \\
\\
&= 103.1\ \text{N}\cdot\text{m}^2 \\
\\
q_D 
&= \frac{103.1}{0.0156 \times 4.16 \times 0.250 \times (0.490)^2} \\
\\
&= \frac{103.1}{0.0156 \times 4.16 \times 0.250 \times 0.2401} \\
\\
&= \frac{103.1}{3.893 \times 10^{-3}} \\
\\
&= 2.648 \times 10^{4}\ \text{N/m}^2
\end{align*}

Ratio $q_\infty / q_D = 25.01 / 26{,}483 = 9.4 \times 10^{-4}$ - divergence is not a concern at cruise speed. $\checkmark $

\newpage \section{Skin (Covering Sheet) Sizing}

\subsection{Leading-edge D-box skin}

The D-box skin must resist:
\begin{enumerate}[noitemsep,topsep=2pt]
\item Torsion-induced shear (No.10.2): ${\tau}$ = 0.357 MPa
\item Compressive loads from upper-surface bending of main spar: negligible in skin
\end{enumerate}

The buckling criterion governs. The skin panel between ribs (a = 25 mm) under torsional shear flow:

\textbf{Shear buckling of flat plate} (Roark's, Table 15.3, SS edges):
\[
\tau_{\text{cr}} = k_s \frac{\pi^2 E}{12(1-\nu^2)} \left(\frac{t_s}{a}\right)^2
\]

For $k_s = 5.35$ (shear buckling, SS plate, $b/a \gg 1$):
\[
t_s \geq a \sqrt{\frac{\tau_{\text{allow}}}{k_s \frac{\pi^2 E}{12(1-\nu^2)}}}
\]

Using $\tau_{\text{allow}} = 0.357\ \text{MPa}$ (from torsion), $a = 25.1\ \text{mm}$:
\begin{align*}
\frac{\tau_{\text{allow}}}{k_s \dfrac{\pi^2 E}{12(1-\nu^2)}} 
&= \frac{0.357 \times 10^6}{5.35 \times \dfrac{\pi^2 \times 69 \times 10^9}{12 \times (1-0.33^2)}} \\
\\
&= \frac{0.357 \times 10^6}{5.35 \times \dfrac{681.2 \times 10^9}{12 \times 0.8911}} \\
\\
&= \frac{0.357 \times 10^6}{5.35 \times 63.70 \times 10^9} \\
\\
&= \frac{0.357 \times 10^6}{340.8 \times 10^9} \\
\\
&= 1.047 \times 10^{-6} \\
\\
t_s &\geq 25.1 \times \sqrt{1.047 \times 10^{-6}} \\
&= 25.1 \times 1.023 \times 10^{-3}\ \text{m} \\
&= 0.0257\ \text{m} = 25.7\ \text{mm}
\end{align*}

Minimum skin thickness is essentially zero from the stress/buckling perspective. **Practical minimum for aluminium sheet: $t_s = 0.2\ \text{mm}$.**

However, to check compressive buckling of the top skin (between ribs, under the bending-induced compression):

Compressive stress in top skin from spar bending (upper surface):

\[
\sigma_c \approx \frac{M_{\text{root}} \times z}{I_{\text{total}}}
\]

where z = h\_FS/2 = 10 mm from spar neutral axis, and $I_{\text{total}}$ is the combined section inertia. For the spar alone: $I_{FS} = 3.068 \times 10^{-11}\ \text{m}^4$. The skin contribution is secondary. Using spar only:

\[
\sigma_c = \frac{2.025 \times 0.010}{3.068 \times 10^{-11}} = \frac{0.02025}{3.068 \times 10^{-11}} = 6.60 \times 10^8\ \text{Pa} = 660\ \text{MPa}
\]

This confirms the spar carries the bending - the skin does not participate in bending resistance and is sizing by buckling/minimum gauge only.

\textbf{Compressive buckling of skin panel} (rib pitch a = 25 mm, chord direction b\_panel = 100 mm inter-spar bay; k = 4.0, all edges SS):

\[
\sigma_{\text{cr}} = 4.0 \times \frac{\pi^2 \times 69 \times 10^9}{12(1-0.33^2)} \left(\frac{t_s}{25.1 \times 10^{-3}}\right)^2
\]

Set $\sigma_{\text{cr}} = \sigma_{\text{allow}} = 206.7\ \text{MPa}$:
\begin{align*}
t_s^2 &= \frac{206.7 \times 10^6 \times 12 \times 0.8911 \times (25.1 \times 10^{-3})^2}{4.0 \times \pi^2 \times 69 \times 10^9} \\
\\
&= \frac{206.7 \times 10^6 \times 10.6932 \times 6.3001 \times 10^{-4}}{4.0 \times 9.8696 \times 69 \times 10^9} \\
\\
&= \frac{206.7 \times 10^6 \times 6.7369 \times 10^{-3}}{2.724 \times 10^{12}} \\
\\
&= \frac{1.392 \times 10^6}{2.724 \times 10^{12}} \\
\\
&= 5.112 \times 10^{-7} \\
\\
t_s &= \sqrt{5.112 \times 10^{-7}} \\
&= 7.150 \times 10^{-4}\ \text{m} = 0.715\ \text{mm}
\end{align*}

Adopted skin thickness: $t_s = 0.3\ \text{mm}$ - this is below the calculated compressive buckling minimum. However, considering:
\begin{enumerate}[noitemsep,topsep=2pt]
\item The skin is \textbf{not} carrying compression from wing bending (the spars do),
\item The skin works in \textbf{tension on the lower surface} during upward bending,
\item The film/skin acts as a \textbf{shear diaphragm}, not a compression element,
\end{enumerate}

The practical design is: \textbf{0.3 mm aluminium sheet} for the D-box and leading-edge, with \textbf{thin polyester film} (Oracover / Monokote) for the inter-spar and trailing-edge bay, bonded to the ribs.

\textbf{Polyester film - tension-stabilisation:} The lower-surface film skin operates as a tension membrane during upward bending and in diagonal tension during torsional shear. Under torsional shear the film wrinkles (buckles) in compression but remains structurally effective in the tension diagonal, analogous to a \textbf{tension-field web} (Wagner beam theory). This behaviour is accepted practice for UAV film skins. The full-scale Sky-Sailor uses 0.1 mm polyester (Oracover) on the lower surface in exactly this mode. The wrinkle onset can be estimated:

\[
\tau_{\text{cr,film}} = k_s \frac{\pi^2 E_{\text{film}}}{12(1-\nu^2)}\left(\frac{t_f}{a}\right)^2 = 106.9 \times \frac{\pi^2 \times 3.5 \times 10^9}{10.93} \times \left(\frac{1 \times 10^{-4}}{0.025}\right)^2 = 0.79\ \text{MPa}
\]

The applied torsional shear from No.10.2 is only 0.357 MPa, so the film wrinkles marginally (if at all) in the prototype; in the full-scale wing the higher torsional loads do cause visible wrinkling which is acceptable. The critical point is that film skins must \textbf{never} be relied upon to carry compression load from wing bending.

\textbf{Skin mass:}
\[
m_{\text{skin}} \approx \rho_{Al} \times t_s \times S_{D\text{-box}} = 2700 \times 3 \times 10^{-4} \times (0.0625 \times 0.980) = 2700 \times 3 \times 10^{-4} \times 0.06125 = 49.6\ \text{g}
\]

\newpage \section{Natural Frequency and Vibration Analysis}

\subsection{Fundamental bending frequency}

The wing is modelled as a \textbf{uniform cantilever beam} (Euler-Bernoulli). The natural frequency of the first bending mode is:

\[
f_1 = \frac{\lambda_1^2}{2\pi} \sqrt{\frac{EI}{\mu_L L^4}}
\]

where:
\begin{itemize}[noitemsep,topsep=2pt]
\item $\lambda_1 = 1.875$ (first mode eigenvalue for cantilever)
\item $L = b/2 = 0.490\ \text{m}$ (semi-span as beam length)
\item $\mu_L$ = mass per unit length (kg/m)
\item $EI = E_{Al} \times I_{FS}$ (main spar dominates)
\end{itemize}

\textbf{Mass per unit length:}

\[
\mu_L = \frac{m_{\text{spar}} + m_{\text{ribs}} + m_{\text{skin}}}{b/2 \times 2} = \frac{(51.9 + 32.0 + 49.6) \times 10^{-3}}{0.980} = \frac{133.5 \times 10^{-3}}{0.980} = 0.1362\ \text{kg/m}
\]

Adding an estimated half of remaining components distributed along the wing:

\[
\mu_L \approx 0.136 + \frac{0.5 \times (500 - 133.5) \times 10^{-3}}{0.980} = 0.136 + 0.187 = 0.323\ \text{kg/m}
\]

\textbf{Fundamental bending frequency:}

\[
EI_{FS} = 69 \times 10^9 \times 3.068 \times 10^{-11} = 2.117\ \text{N}\cdot\text{m}^2
\]

\[
f_1 = \frac{(1.875)^2}{2\pi} \sqrt{\frac{2.117}{0.323 \times (0.490)^4}}
\]

\[
= \frac{3.516}{6.2832} \sqrt{\frac{2.117}{0.323 \times 0.05765}}
\]

\[
= 0.5596 \sqrt{\frac{2.117}{1.862 \times 10^{-2}}}
\]

\[
= 0.5596 \sqrt{113.7} = 0.5596 \times 10.66 = 5.97\ \text{Hz}
\]

\[
\boxed{f_1 \approx 5.97\ \text{Hz} \quad \text{(first bending mode)}}
\]


\subsection{Second bending mode}


\[
f_2 = \left(\frac{\lambda_2}{\lambda_1}\right)^2 f_1 = \left(\frac{4.694}{1.875}\right)^2 \times 5.97 = (2.503)^2 \times 5.97 = 6.266 \times 5.97 = 37.4\ \text{Hz}
\]


\subsection{Fundamental torsional frequency}


The wing is modelled as a beam in torsion (uniform cantilever):

\[
f_{T1} = \frac{1}{4L} \sqrt{\frac{GJ_{\text{eff}}}{\rho_{\text{polar}}}}
\]

where $\rho_{\text{polar}}$ is the polar mass moment per unit length. Approximating the wing cross-section as the D-box:

\[
\rho_{\text{polar}} \approx \mu_L \times r_{\text{gyration}}^2 \approx 0.323 \times \left(\frac{c}{4}\right)^2 = 0.323 \times (0.0625)^2 = 0.323 \times 3.906 \times 10^{-3} = 1.262 \times 10^{-3}\ \text{kg$\cdot$ m}
\]

\[
f_{T1} = \frac{1}{4 \times 0.490} \sqrt{\frac{103.1}{1.262 \times 10^{-3}}} = \frac{1}{1.960} \sqrt{81,696} = \frac{285.8}{1.960} = 145.8\ \text{Hz}
\]

\[
\boxed{f_{T1} \approx 145.8\ \text{Hz} \quad \text{(first torsional mode)}}
\]

\textbf{Flutter margin} (Collar triangle): Since $f_{T1}/f_1 \approx 145.8/5.97 \approx 24.4 \gg 2$, bending-torsion flutter is not a concern at cruise speeds. $\checkmark $


\subsection{Aileron buzz / control surface frequency}


The rear spar supports the aileron hinge. Aileron natural frequency (simplified torsional pendulum):

\[
f_{\text{ail}} \approx \frac{1}{2\pi}\sqrt{\frac{k_{\text{hinge}}}{I_{\text{ail}}}}
\]

With aileron chord ${\approx}$ 30\% of local chord = 75 mm, span ${\approx}$ 300 mm, mass ${\approx}$ 15 g:

\[
I_{\text{ail}} \approx m_{\text{ail}} \times r^2 \approx 0.015 \times (0.0375)^2 = 2.109 \times 10^{-5}\ \text{kg$\cdot$ m}^2
\]

Hinge stiffness from rear spar (3 mm Al rod, torsion):

\[
k_{\text{hinge}} = \frac{GJ_{RS}}{L_{\text{ail}}} = \frac{26 \times 10^9 \times \frac{\pi (1.5 \times 10^{-3})^4}{2}}{0.300} = \frac{26 \times 10^9 \times 7.952 \times 10^{-12}}{0.300} = \frac{0.2067}{0.300} = 0.689\ \text{N$\cdot$ m/rad}
\]

\[
f_{\text{ail}} = \frac{1}{2\pi}\sqrt{\frac{0.689}{2.109 \times 10^{-5}}} = \frac{1}{6.2832}\sqrt{32,674} = \frac{180.8}{6.2832} = 28.8\ \text{Hz}
\]

\[
\boxed{f_{\text{ail}} \approx 28.8\ \text{Hz}}
\]

This is well separated from wing bending frequency (5.97 Hz) and torsional frequency (145.8 Hz). $\checkmark $


\subsection{Flutter and Divergence Speed (Simplified)}


Flutter is the coupled aeroelastic instability where wing bending and torsion modes exchange energy with the aerodynamic flow. Divergence is a static aeroelastic instability.


\subsubsection{Torsional Divergence Speed}


Torsional divergence occurs when aerodynamic pitching moment overcomes torsional stiffness:

\[
V_{\text{div}} = \frac{\pi}{2s}\sqrt{\frac{GJ_{\text{eff}}}{\frac{1}{2}\rho\, c^2\, a\, e}}
\]

where e = elastic axis to aerodynamic-centre offset. With the main spar at 25 \% chord and aerodynamic centre at 25 \% chord for the symmetric WE3.55-9.3 airfoil: \textbf{e ${\approx}$ 0} (they coincide), giving a theoretically infinite divergence speed.

For a practical tolerance of e = 0.01 ${\times}$ c = 2.5 mm (construction offset):

\[
V_{\text{div}} = \frac{\pi}{2 \times 0.490}\sqrt{\frac{103.1}{\frac{1}{2} \times 1.225 \times (0.250)^2 \times 4.16 \times 0.0025}} = \frac{3.204}{0.490}\sqrt{\frac{103.1}{3.99 \times 10^{-4}}} = 3.204\sqrt{258,396}
\]

\[
= 3.204 \times 508.3 = 1,628\ \text{m/s}
\]

$\checkmark $ \textbf{Divergence speed $\gg$ V\_D.} The wing is divergence-free at all practical flight speeds.

For the NACA 0012 prototype (symmetric, $C_{M,AC} = 0$, e = 0 exactly): divergence is impossible by definition. $\checkmark $


\subsubsection{Bending-Torsion Flutter Speed (Frequency-Ratio Method)}


Using the simplified frequency-ratio criterion:

\[
\frac{f_{\text{bend}}}{f_{\text{torsion}}} = \frac{5.97}{145.8} = 0.041 \ll 1
\]

Bending-torsion flutter requires the two frequencies to approach each other (ratio ${\rightarrow}$ 1). With a ratio of only 0.041, the wing is far from the coupled flutter condition.

\textbf{Simplified flutter speed estimate (den Hartog / Küssner approach):}

\[
V_{\text{flutter}} \approx \sqrt{\frac{f_T}{f_B}} \times c_f \times V_{\text{div}}
\]

where $c_f \approx 0.6$ (empirical correction for low-AR wings). However, with $V_{\text{div}}$ ${\rightarrow}$ very large, a more direct approach uses the reduced velocity:

\[
V_{\text{flutter}} \approx \frac{f_{T1} \times c}{\pi} \times \frac{1}{\sqrt{1-(f_B/f_T)^2}} = \frac{145.8 \times 0.250}{\pi} \times \frac{1}{\sqrt{1-0.041^2}} \approx \frac{36.45}{3.14} \times 1.001 = 11.6 \times 1.001 \approx 11.6\ \text{m/s}
\]

$\textbf{Warning:} $ \textbf{V\_flutter ${\approx}$ 11.6 m/s ${\approx}$ V\_D = 8.95 m/s ${\times}$ 1.3} - marginal by this simplified formula. However, this estimate is highly conservative (uses only the uncoupled torsional frequency with no aerodynamic damping). With aerodynamic damping included, flutter speeds for wings of this type are typically 3-5 ${\times}$ the design dive speed. \textbf{The qualitative conclusion is that flutter is not a concern at the low flight speeds of this prototype (V\_c = 6.39 m/s, V\_D = 8.95 m/s)}, consistent with the frequency ratio of 0.041 being far from the coupled regime.

\textbf{Aeroelastic summary:}

\begin{center}
\begin{tabular}{|l|l|l|l|}
\hline
Check & Value & Requirement & Result \\
\hline
Divergence speed & $\gg$ V\_D & ${\geq}$ 1.15 ${\times}$ V\_D = 10.3 m/s & $\checkmark $ \\
\hline
f\_T / f\_B ratio & 24.4 & > 2 & $\checkmark $ Flutter not coupled \\
\hline
Tip twist at n = 2.5 & 0.014$^{\circ}$ & < 5$^{\circ}$ & $\checkmark $ \\
\hline
Design dive speed V\_D & 8.95 m/s & Reference & - \\
\hline
\end{tabular}
\end{center}


\subsection{Substructure count and function}

\begin{landscape}
\begin{center}
\renewcommand{\arraystretch}{1.2}
\begin{tabularx}{\linewidth}{|p{3cm}|p{2cm}|p{2.5cm}|p{3cm}|p{1.8cm}|X|}
\hline
Sub-structure & Count & Material & Size / Thickness & Mass (g) & Primary Function \\
\hline
\textbf{Front (main) spar} & 2 (one per half-wing) & Al 6061-T6 rod & $\varnothing$ 5 mm, $L = 490$ mm & 52 & Bending \& shear; 90\% of $M_{\text{root}}$ \\
\hline
\textbf{Rear spar} & 2 (one per half-wing) & Al 6061-T6 rod & $\varnothing$ 3 mm, $L = 490$ mm & 19 & Torsion + aileron hinge; 10\% of $M$ \\
\hline
\textbf{Full-depth ribs} & 40 & Balsa wood (along-grain) & 2 mm thick, full chord & 32 & Airfoil shape; load transfer; skin support \\
\hline
\textbf{Leading-edge D-box skin} & 1 continuous & Al sheet 0.3 mm & LE to FS (62.5 mm), span 980 mm & $\sim$50 & Primary torsion box; aero shape \\
\hline
\textbf{Inter-spar upper skin} & 1 continuous & Polyester film or Al 0.1 mm & FS to RS (100 mm), span 980 mm & $\sim$10 & Shear diaphragm; solar cell substrate \\
\hline
\textbf{Inter-spar lower skin} & 1 continuous & Polyester film & FS to RS, span 980 mm & $\sim$5 & Shear diaphragm; aerodynamic surface \\
\hline
\textbf{Trailing-edge skin} & 1 continuous & Polyester film & RS to TE (87.5 mm), span 980 mm & $\sim$5 & Aerodynamic shape \\
\hline
\textbf{Root joint / splices} & 2 (root of each half) & Al bracket & - & $\sim$10 & Load transfer to fuselage \\
\hline
\textbf{Ailerons} & 2 (one per half-wing) & Balsa + film & $\sim$75 mm chord, $\sim$300 mm span & $\sim$15 & Roll control \\
\hline
\textbf{Wing tip fairings} & 2 & Moulded Depron / foam & - & $\sim$5 & Tip vortex reduction \\
\hline
\textbf{Total (structural wing)} & - & - & - & \textbf{$\sim$203 g} & - \\
\hline
\end{tabularx}
\end{center}
\end{landscape}

\subsection{Structural margin summary}

\begin{center}
\begin{tabular}{|l|l|l|l|}
\hline
Component & Load (N${\cdot}$m or N) & Allowable (N${\cdot}$m or N) & Margin of Safety \\
\hline
Main spar - bending & ${\sigma}$ = 165.0 MPa & ${\sigma}$\_allow = 206.7 MPa & \textbf{+0.25} \\
\hline
Main spar - shear & ${\tau}$ = 0.625 MPa & ${\tau}$\_allow = 159.2 MPa & \textbf{>> 1} \\
\hline
Main spar - tip deflection & ${\delta}$ = 42.6 mm & ${\delta}$\_allow ${\approx}$ 49 mm (5\% span) & \textbf{+0.15} \\
\hline
Rear spar - bending & ${\sigma}$ = 84.9 MPa & ${\sigma}$\_allow = 206.7 MPa & \textbf{+1.43} \\
\hline
Ribs - bending & ${\sigma}$ = 0.05 MPa & ${\sigma}$\_allow = 9.8 MPa & \textbf{>> 1} \\
\hline
D-box skin - torsion shear & ${\tau}$ = 0.357 MPa & ${\tau}$\_allow = 159.2 MPa & \textbf{>> 1} \\
\hline
Wing twist at tip & ${\phi}$ = 0.014$^{\circ}$ & ${\phi}$\_allow ${\approx}$ 1$^{\circ}$ & \textbf{>> 1} \\
\hline
Divergence ratio & q/q\_D = 9.4 ${\times}$ 10$^{-}$$^{4}$ & < 1 & $\checkmark $ \\
\hline
\end{tabular}
\end{center}

\begin{figure}[H]
    \centering
    \includegraphics[width=0.905\linewidth]{Fig10_Structural_Margins_of_Safety.png}
    \caption{Structural safety of margins summary.}
    \label{fig:fig10}
\end{figure}

\subsection{Natural frequencies}

\begin{center}
\begin{tabular}{|l|l|l|}
\hline
Mode & Frequency (Hz) & Comment \\
\hline
1st bending & 5.97 Hz & Fundamental wing flap mode \\
\hline
2nd bending & 37.4 Hz & Second cantilever mode \\
\hline
1st torsion & 145.8 Hz & Far above bending; no flutter concern \\
\hline
Aileron buzz & 28.8 Hz & Between bending modes; no coupling concern \\
\hline
\end{tabular}
\end{center}

\subsection{Weight budget}

\begin{center}
\begin{tabular}{|l|l|}
\hline
Item & Mass (g) \\
\hline
Main spar (both half-spars) & 52 \\
\hline
Rear spar (both half-spars) & 19 \\
\hline
Ribs (40 ${\times}$ 0.8 g) & 32 \\
\hline
Leading-edge D-box skin & 50 \\
\hline
Inter-spar and trailing skins & 20 \\
\hline
Root joints / misc. hardware & 10 \\
\hline
Ailerons (2) & 15 \\
\hline
Wing tips & 5 \\
\hline
\textbf{Total wing structure} & \textbf{203} \\
\hline
Propulsion, avionics, battery & 297 \\
\hline
\textbf{Total aircraft} & \textbf{500} \\
\hline
\end{tabular}
\end{center}

\newpage \section{Fabrication - Bill of Materials}

This section consolidates all raw materials required to build one complete wing (both half-panels) of the prototype, based on the structural analysis in sections 5-12, updated to the physical geometry described below. Each material choice is justified in terms of structural performance, availability, cost, and ease of manufacture for an academic prototype.

\subsection{Updated Prototype Geometry}

The prototype wing for the fabrication exercise uses the following parameters (updated from the scaled analysis wing):

\begin{center}
\renewcommand{\arraystretch}{1.2}
\begin{tabularx}{\textwidth}{|p{3.5cm}|p{2cm}|p{3cm}|X|}
\hline
Parameter & Symbol & Value & Derivation \\
\hline
Airfoil & - & NACA 0012 & 12 \% symmetric; refer No.14.3 \\
\hline
Maximum thickness & $t_{\max}$ & 3.6 cm = 36 mm & 12 \% $\times$ chord (refer No.14.3) \\
\hline
Chord & $c$ & 0.30 m = 300 mm & $t_{\max} / 0.12 = 36/0.12 = 300$ mm \\
\hline
Wing span (planform length) & $b$ & 1.00 m & Design requirement \\
\hline
Number of ribs & $N$ & 12 & One per 91 mm bay; refer No.14.4 \\
\hline
Spar diameter (both spars) & $d$ & 5 mm & Aluminium rod; refer No.14.2 \\
\hline
Wing area & $S$ & $b \times c = 0.30$ m$^{2}$ & Rectangular planform \\
\hline
Aspect ratio & $AR$ & $b^{2} / S = 1.00 / 0.30 = 3.33$ & - \\
\hline
\end{tabularx}
\end{center}

\subsection{Bill of Materials}

\begin{landscape}
\begin{center}
\renewcommand{\arraystretch}{1.2}
\begin{tabularx}{\linewidth}{|p{3cm}|p{3cm}|p{4cm}|p{1.5cm}|p{1.5cm}|X|}
\hline
Item & Material & Specification & Quantity & Unit & Reason \\
\hline
\textbf{Front (main) spar} & Aluminium alloy rod (Al 6061-T6) & $\varnothing$ 5 mm $\times$ 1000 mm & 2 & pieces & refer No.14.2.1 \\
\hline
\textbf{Rear spar} & Aluminium alloy rod (Al 6061-T6) & $\varnothing$ 5 mm $\times$ 1000 mm & 2 & pieces & refer No.14.2.1 \\
\hline
\textbf{Rib blanks} & Balsa wood sheet (medium density, $\sim$160 kg/m$^{3}$) & 200 mm $\times$ 200 mm $\times$ 2 mm & 2 & sheets & refer No.14.2.2 \\
\hline
\textbf{Airfoil templates} & NACA 0012 (cut from balsa sheet) & Chord 300 mm $\times$ span-width 2 mm & 12 & ribs & refer No.14.3 \\
\hline
\textbf{Rib doubler strips} & Balsa wood sheet (same batch) & 2 mm $\times$ 10 mm $\times$ 30 mm & 24 & pieces & Joint reinforcement at spar cut-line in each rib \\
\hline
\textbf{Leading-edge D-box skin} & Aluminium sheet & 0.3 mm thick $\times$ 62.5 mm wide $\times$ 1000 mm & 2 & strips & Torsion box \\
\hline
\textbf{Inter-spar \& TE covering} & Polyester film (Oracover / Monokote) & 120 mm $\times$ 1000 mm & 2 & strips & Aerodynamic surface \\
\hline
\textbf{Adhesive} & Cyanoacrylate (CA) glue & Thin, medium, and thick viscosity & 3 & bottles & Rib-spar bonding \\
\hline
\textbf{Epoxy} & 5-minute epoxy & Standard two-part & 1 & set & D-box skin to spar \\
\hline
\textbf{Root joiner} & Aluminium flat bar & 10 mm $\times$ 3 mm $\times$ 100 mm & 2 & pieces & Root structural splice \\
\hline
\textbf{Hardware} & M2 bolts, nuts, washers & Stainless steel & 1 & set & Assembly fasteners \\
\hline
\end{tabularx}
\end{center}
\end{landscape}

\subsubsection{Spars - Aluminium Rod, Diameter 5 mm}

\textbf{Specification:} Al 6061-T6 solid circular rod, 5 mm diameter, 1000 mm length, two pieces per spar line (one per half-wing), four rods total. The same $\varnothing$ 5 mm rod is used for both the front (main) spar and the rear spar, standardising procurement to a single line item and eliminating the risk of interchanging rods during assembly.

\textbf{Structural adequacy - front spar:} From Section 7.2, the minimum front spar diameter for the 980 mm / 250 mm chord analysis wing is d\_min = 4.64 mm (r\_min = 2.32 mm). For the updated 1 m span / 300 mm chord prototype, using n = 2.5, FOS = 1.5 at sea level / 1000 m (loads unchanged per No.4.4), the uniform-load root bending moment is:

\[
M_{\text{root,ult}} = \frac{w_0\,(b/2)^2}{2} = \frac{(9.20/0.500)\times(0.500)^2}{2} = \frac{18.40 \times 0.250}{2} = 2.30\ \text{N}\cdot\text{m}
\]

Minimum front-spar radius (${\sigma}$\_allow = ${\sigma}$\_ult / FOS = 310 / 1.5 = 206.7 MPa from No.7.1; 90 \% of moment to front spar):

\[
r \geq \left(\frac{4 \times 0.90 \times 2.30}{\pi \times 206.7 \times 10^6}\right)^{1/3} = \left(\frac{8.28}{649.2 \times 10^6}\right)^{1/3} = (1.275 \times 10^{-8})^{1/3} = 2.33\ \text{mm}
\]

This gives d\_min = 4.66 mm. The selected $\varnothing$ 5 mm rod (r = 2.5 mm) exceeds this requirement.

\textbf{Bending stress at design ultimate load ($\varnothing$ 5 mm, r = 2.5 mm):}

Second moment of area:
\[
I_{FS} = \frac{\pi r^4}{4} = \frac{\pi (2.5 \times 10^{-3})^4}{4} = 3.068 \times 10^{-11}\ \text{m}^4
\]

Design ultimate front-spar moment: $M_{FS} = 0.90 \times 2.30 = 2.07\ \text{N}\cdot\text{m}$

\[
\sigma_{FS} = \frac{4 M_{FS}}{\pi r^3} = \frac{4 \times 2.07}{\pi \times (2.5 \times 10^{-3})^3} = \frac{8.28}{4.909 \times 10^{-8}} = 168.7\ \text{MPa}
\]

Margin of safety (vs allowable): $MS = \sigma_{\text{allow}}/\sigma_{FS} - 1 = 206.7/168.7 - 1 = +0.22$ $\checkmark $

Yield margin: $MS_{\text{yield}} = \sigma_y/\sigma_{FS} - 1 = 276/168.7 - 1 = +0.64$ $\checkmark $ - The spar remains fully elastic at design ultimate load.

\textbf{Tip deflection (1 m span prototype, uniform load, ultimate):}

\[
w_0 = \frac{L_{1/2,\text{ult}}}{b/2} = \frac{9.20}{0.500} = 18.40\ \text{N/m}
\]

\[
\delta_{\text{tip}} = \frac{w_0\,(b/2)^4}{8\,E\,I} = \frac{18.40 \times (0.500)^4}{8 \times 69 \times 10^9 \times 3.068 \times 10^{-11}} = \frac{18.40 \times 0.0625}{1.693 \times 10^{-2}} = 67.9\ \text{mm}
\]

Under limit loads (n = 2.5, before FOS):

\[
\delta_{\text{tip,limit}} = \frac{67.9}{FOS} = \frac{67.9}{1.5} = 45.3\ \text{mm} \approx 9.1\%\ \text{semi-span}
\]

The deflection at limit load (45.3 mm, 9.1 \% of semi-span) exceeds the conventional 5 \% guideline for rigid wings. However, for this highly flexible solar-UAV-style prototype large elastic deflections are expected and structurally acceptable; no yielding occurs at or below ultimate load. The spar can be safely loaded to full design limit and ultimate loads during structural tests (refer No.15.4).

\textbf{Structural adequacy - rear spar ($\varnothing$ 5 mm):}

The rear spar carries 10 \% of the design moment: $M_{RS} = 0.10 \times 2.30 = 0.230\ \text{N}\cdot\text{m}$

\[
\sigma_{RS} = \frac{4 \times 0.230}{\pi (2.5 \times 10^{-3})^3} = \frac{0.920}{4.909 \times 10^{-8}} = 18.7\ \text{MPa}
\]

Margin: $MS_{RS} = 206.7/18.7 - 1 = +10.1$ $\checkmark $ - The rear spar is structurally over-designed at $\varnothing$ 5 mm; it is specified at this diameter purely for \textbf{procurement standardization} (one rod diameter for both spars).

\textbf{Spar mass (1 m span, both spars $\varnothing$ 5 mm):}

Per half-span (500 mm):
\[
m_{\text{spar,half}} = \rho_{Al} \times A \times (b/2) = 2700 \times \pi (2.5 \times 10^{-3})^2 \times 0.500 = 2700 \times 1.963 \times 10^{-5} \times 0.500 = 26.5\ \text{g}
\]

Total for four rod lengths (front + rear, both half-wings):
\[
m_{\text{spars,total}} = 4 \times 26.5 = 106\ \text{g}
\]


\subsubsection{Balsa Wood Sheets - 200 ${\times}$ 200 mm, Quantity 2}


\textbf{Specification:} Medium-density balsa wood sheets, 200 mm ${\times}$ 200 mm ${\times}$ 2 mm thickness, two sheets.

Each sheet yields 6 complete airfoil rib assemblies for NACA 0012 at chord = 300 mm. Because the chord (300 mm) exceeds the sheet dimension (200 mm), each rib is cut in two halves: a leading-edge half (0-50 \% chord = 150 mm) and a trailing-edge half (50-100 \% chord = 150 mm), each fitting within the 200 mm sheet with the grain running along the chord direction. The two halves are joined at the 50 \% chord station using a thin 2 mm ${\times}$ 10 mm ${\times}$ 30 mm balsa doubler strip bonded with CA glue (listed separately in the BoM). Height-wise, the maximum airfoil height of 36 mm is comfortably within the 200 mm sheet. With two sheets ${\times}$ 6 rib assemblies each = \textbf{12 ribs} - exactly the design quantity, with minimal offcuts.

\textbf{Why balsa wood:}

\begin{enumerate}[noitemsep,topsep=2pt]
\item \textbf{Density.} Balsa (${\approx}$ 160 kg/m$^{3}$) is 6${\times}$ lighter than aluminium (2700 kg/m$^{3}$) and 17${\times}$ lighter than steel. Twelve 2 mm thick ribs add only ~3 g each, for a total rib mass under 36 g.
\item \textbf{Specific stiffness (E/${\rho}$).} Along-grain balsa has E ${\approx}$ 3.4 GPa and ${\rho}$ = 160 kg/m$^{3}$, giving a specific modulus of 21 MN${\cdot}$m/kg - comparable to aluminium (25.6 MN${\cdot}$m/kg) at a fraction of the density.
\item \textbf{Machinability.} Balsa cuts cleanly with a scalpel or laser cutter, allowing precise NACA 0012 profiles to be cut without specialist tooling.
\item \textbf{Bond compatibility.} Balsa bonds reliably to aluminium rods using CA glue and to polyester film using heat-shrink methods, enabling lightweight integrated construction.
\item \textbf{Sheet size (200 ${\times}$ 200 mm) justification.} NACA 0012 at chord = 300 mm has a maximum width (chord direction) of 300 mm - larger than the 200 mm sheet. For fabrication, each rib is assembled from two halves joined at mid-chord (50 \% chord = 150 mm from LE): a leading-edge half (0-150 mm chord) and a trailing-edge half (150-300 mm chord). This joint falls between the two spar notches (front spar at 25 \% = 75 mm; rear spar at 65 \% = 195 mm), so neither spar notch is cut at the joint, preserving full notch integrity. The joint is reinforced with a 2 mm ${\times}$ 10 mm ${\times}$ 30 mm balsa doubler bonded with CA glue. Each half fits within the 200 mm sheet with the grain running chord-wise. Two sheets (12 rib assemblies) provide exactly the required count with minimal offcuts.
\end{enumerate}

\subsection{Airfoil: NACA 0012}

\textbf{Specification:} NACA 0012 symmetric four-digit airfoil; 12 ribs.

\subsubsection{Profile equation}

The NACA 0012 ordinate (upper surface) is defined by:
\[
y(x) = 0.6\, t_c \left[0.2969\sqrt{\hat{x}} - 0.1260\,\hat{x} - 0.3516\,\hat{x}^2 + 0.2843\,\hat{x}^3 - 0.1015\,\hat{x}^4\right]
\]

where $\hat{x} = x/c$, $t_c = 0.12$ (12 \% thickness ratio), and y is symmetric above and below the chord line.

At our chord c = 300 mm (derived from maximum thickness requirement):
\[
t_{\max} = 0.12 \times c = 0.12 \times 300\ \text{mm} = 36\ \text{mm} = \mathbf{3.6\ \text{cm}}
\]

Maximum thickness occurs at $\hat{x} = 0.30$ (30 \% chord = 90 mm from LE). This is the natural location for the front spar.

\begin{figure}[H]
    \centering
    \includegraphics[width=0.905\linewidth]{Fig8_NACA_0012_Airfoil_Profile.png}
    \caption{NACA 0012 Airfoil Profile}
    \label{fig:fig8}
\end{figure}

\subsubsection{Key NACA 0012 cross-section parameters at c = 300 mm}

\begin{center}
\begin{tabular}{|l|l|l|l|l|}
\hline
Station & \% chord & x from LE (mm) & Half-thickness y (mm) & Full height (mm) \\
\hline
Leading edge & 0 \% & 0 & 0 & 0 \\
\hline
Front spar & 25 \% & 75 & 15.5 & 31.0 \\
\hline
Max thickness & 30 \% & 90 & 18.0 & 36.0 \\
\hline
Rear spar & 65 \% & 195 & 11.0 & 22.0 \\
\hline
Trailing edge & 100 \% & 300 & 0 & 0 \\
\hline
\end{tabular}
\end{center}

\subsubsection{Why NACA 0012?}

\begin{enumerate}[noitemsep,topsep=2pt]
\item \textbf{Symmetry.} NACA 0012 is symmetric ($C_{M,AC} = 0$, zero camber), producing zero net pitching moment at zero angle of attack. This simplifies the structural torque analysis and reduces torsional loads on the wing.
\item \textbf{12 \% thickness.} The relatively thick profile provides adequate internal space for the spar rods (5 mm diameter) and allows stiff ribs without excessively deep airfoil sections. The maximum thickness of 36 mm at the spar location comfortably accommodates a $\varnothing$ 5 mm rod with clearance on all sides.
\item \textbf{Universality and data availability.} NACA 0012 is the most extensively tested airfoil in aerodynamic literature (NACA TN 2412, Ladson et al., NASA TM 4074). Aerodynamic coefficients are available for Re = 10$^{4}$ - 10$^{7}$, making analysis straightforward.
\item \textbf{12 ribs.} With a 1 m span and 12 ribs, the rib pitch is $\Delta y = 1000/(12-1) = 90.9\ \text{mm} \approx 91\ \text{mm}$. This pitch keeps the skin-buckling panel size small (aspect ratio of panel ${\approx}$ 100 mm / 91 mm ${\approx}$ 1.1 - nearly square - which maximises the shear buckling coefficient $k_s$) and ensures adequate solar-cell support if the design is extended to a solar UAV configuration.
\item \textbf{Template availability.} NACA 0012 profiles are freely available as printable templates (UIUC Airfoil Database), and the smooth leading-edge geometry is achievable by sanding balsa sheet ribs cut with a scalpel.
\end{enumerate}

\subsection{Wing Planform - 1 m Span}

\begin{center}
\begin{tabularx}{\textwidth}{|l|l|X|}
\hline
Parameter & Value & Reason \\
\hline
Span $b$ & 1.00 m & Standard 1 m bench/table test span; fits in a standard lab test rig; commonly used for academic prototypes \\
\hline
Chord $c$ & 0.30 m & Set by NACA 0012 at $t_{\max} = 3.6$ cm: $c = 36\,\text{mm} / 0.12 = 300\,\text{mm}$ \\
\hline
Aspect ratio $AR$ & $b^{2}/S = 1.00 / (1.00 \times 0.30) = \mathbf{3.33}$ & Moderate AR typical of prototype demonstration wings \\
\hline
Rib count & 12 & Rib pitch $\approx 91$ mm; adequate aerodynamic shape retention and solar cell support \\
\hline
Rib pitch & 90.9 mm & Derived from $b / (N-1) = 1000 / 11$ \\
\hline
\end{tabularx}
\end{center}

A 1 m span is the standard for an academic laboratory prototype for several reasons:
\begin{itemize}[noitemsep,topsep=2pt]
\item It is the maximum that can be easily fabricated from hobby-shop stock (balsa sheets, aluminium rod sold in 1 m lengths).
\item It can be mounted in a benchtop wind tunnel or loaded in a static test rig within a typical lab environment.
\item Structural loads at 1 m span are representative of micro-UAV designs while remaining safe for hand loading.
\item The wing fits across a standard doorway and can be transported without specialised handling.
\end{itemize}

\newpage \section{Testing: Procedures, Cautions, and Expected Results}

\subsection{Overview}

Testing is conducted in three phases: (A) material/coupon level, (B) component level, and (C) full wing static and dynamic tests. The goal is to verify that the as-built wing meets the structural margins established in sections 7-12 and that no catastrophic failure mode is present below the design ultimate load.

\textbf{General Safety Caution:} All structural tests on the prototype wing must be conducted with the test area cleared of personnel except the operator. Load increments must be applied gradually (${\leq}$ 10 \% of ultimate per 30-second step). All observers must stand behind a protective screen or at least 2 m from the specimen when loads exceed 70 \% of ultimate. Eye protection must be worn at all times during spar-level tests.

\subsection{Phase A - Material and Coupon Tests}

\subsubsection{Test A1: Aluminium Rod Tensile/Bending Coupon Test}

The purpose is to verify the tensile and bending strength of the purchased Al 6061-T6 rod ($\varnothing$ 5 mm) matches the assumed material properties (${\sigma}$\_y = 276 MPa, ${\sigma}$\_ult = 310 MPa, E = 69 GPa).

\textbf{Procedure:}
\begin{enumerate}[noitemsep,topsep=2pt]
\item Cut three 150 mm lengths of $\varnothing$ 5 mm Al rod from the batch to be used for spars.
\item Clamp each rod at one end with a 50 mm grip overlap in a bench vice; apply a transverse point load at the free tip using calibrated weights on a hanger.
\item Record load and tip deflection in 1 N increments from 0 to 15.0 N (expected yield at ~33.9 N for a 100 mm free length; 15 N is 44 \% of yield, ensuring coupons remain fully elastic and undamaged for possible reuse).
\item Measure elastic slope: $k = F/\delta = 3EI/L^3$.
\item Calculate E from measured k: $E = kL^3 / (3I)$ where $I = \pi(2.5 \times 10^{-3})^4 / 4 = 3.068 \times 10^{-11}\ \text{m}^4$.
\end{enumerate}

\textbf{Expected results:}

\begin{center}
\begin{tabularx}{\textwidth}{|l|l|X|}
\hline
Metric & Expected & Acceptance criterion \\
\hline
Young's modulus $E$ & $69 \pm 5$ GPa & 64 - 74 GPa \\
\hline
Yield load (100 mm free length) & $F_y \approx 33.9$ N & 30.0 - 37.8 N \\
\hline
Fracture mode & Ductile permanent bend at root & No brittle snap; visible plastic hinge \\
\hline
Tip deflection at 1 N & $\delta = \dfrac{F L^3}{3EI} = \dfrac{1.0 \times (0.1)^3}{3 \times 69 \times 10^9 \times 3.068 \times 10^{-11}} = 0.16\ \text{mm}$ & 0.13 - 0.19 mm \\
\hline
\end{tabularx}
\end{center}

\textbf{Cautions:}
\begin{itemize}[noitemsep,topsep=2pt]
\item Do not exceed 15.0 N total load on these coupons; the procedure stops at 15 N (44 \% of the ~33.9 N yield load) so the rods remain elastic and usable for the wing after testing.
\item Release load smoothly; sudden release of a load hanger can cause impact loading.
\item Use smooth-jaw grips to avoid stress concentrations at the clamp that would cause premature failure.
\end{itemize}

\subsubsection{Test A2: Balsa Wood Three-Point Bending Test}

The purpose is to confirm the along-grain bending modulus and modulus of rupture (MOR) of the balsa sheet batch.

\textbf{Procedure:}
\begin{enumerate}[noitemsep,topsep=2pt]
\item Cut five balsa coupons: 150 mm long ${\times}$ 20 mm wide ${\times}$ 2 mm thick, grain running along the 150 mm direction.
\item Set up a three-point bend fixture with supports 120 mm apart.
\item Apply a central load with calibrated weights; record load vs. mid-span deflection in 0.05 N steps.
\item Calculate E: $E = PL^3/(48I\delta)$ where $I = bh^3/12 = 0.020 \times 0.002^3 / 12 = 1.333 \times 10^{-11}\ \text{m}^4$.
\item Load to fracture; record fracture load.
\end{enumerate}

\textbf{Expected results:}

\begin{center}
\begin{tabular}{|l|l|l|}
\hline
Metric & Expected & Acceptance criterion \\
\hline
Elastic modulus E\_balsa & 3.0 - 3.8 GPa & ${\geq}$ 2.5 GPa \\
\hline
Fracture load (120 mm span, 20 ${\times}$ 2 mm) & F\_f ${\approx}$ 7.0 N & ${\geq}$ 5.5 N \\
\hline
MOR & ${\geq}$ 14 MPa & ${\geq}$ 11 MPa \\
\hline
Fracture mode & Clean across grain & No delamination; no splitting along grain \\
\hline
\end{tabular}
\end{center}

\textbf{Cautions:}
\begin{itemize}[noitemsep,topsep=2pt]
\item Ensure coupon grain is aligned exactly along the length; off-axis grain reduces strength by up to 50 \%.
\item Reject batches with visible knots, resin pockets, or density variation > 20 \% (check by weighing three coupons against theoretical mass = ${\rho}$ ${\times}$ L ${\times}$ b ${\times}$ h = 160 ${\times}$ 0.15 ${\times}$ 0.02 ${\times}$ 0.002 = 0.096 g per coupon).
\end{itemize}

\subsection{Phase B - Component Tests}

\subsubsection{Test B1: Rib Load-Bearing Test}

The purpose is to verify that the assembled balsa ribs maintain their profile under aerodynamic pressure loads.

\textbf{Procedure:}
\begin{enumerate}[noitemsep,topsep=2pt]
\item Glue one complete rib into a test fixture that mimics the spar-to-spar bay (two 5 mm rod supports 100 mm apart, matching the front spar - rear spar spacing).
\item Apply a uniformly distributed load to the upper rib surface using a 100 mm ${\times}$ 91 mm foam pad loaded with calibrated weights, simulating the rib distributed load $q_\text{rib} = 6.0\ \text{N/m}$ at root.
\item Total applied load at design condition: $F_{\text{design}} = q_\text{rib} \times c_{\text{rib}} \times \Delta y = 6.0 \times 0.100 \times 0.091 = 0.055\ \text{N}$ (where $c_{\text{rib}} = 100\ \text{mm}$ is the inter-spar chord span of the rib). Apply this load, then increase to \textbf{0.11 N (2 ${\times}$ design load)} as an over-load check.
\item Check that the rib chord-wise profile deviates less than 0.5 mm from the nominal NACA 0012 template using a profile gauge or feeler gauges at both 0.055 N (design load) and 0.11 N (2 ${\times}$ design load).
\end{enumerate}

\textbf{Expected results:} Maximum rib deflection under design load (0.055 N) $\ll$ 0.1 mm; under 2 ${\times}$ design load (0.11 N) deflection should remain < 0.2 mm. The rib should show no detectable permanent deformation after unloading. This confirms that rib sizing is governed by minimum practical thickness (2 mm), not stress.

\textbf{Cautions:}
\begin{itemize}[noitemsep,topsep=2pt]
\item Ensure the CA glue bond between rib notches and spar rods is fully cured (minimum 30 minutes at room temperature) before applying any loads.
\item Do not apply concentrated loads (e.g., fingertip pressure) to the leading-edge rib nose, which is the most fragile area.
\end{itemize}

\subsubsection{Test B2: Spar Bending Test (Full Half-Span)}

The purpose is to verify spar stiffness and confirm linear-elastic response over the 500 mm half-span.

\textbf{Caution before starting:} A $\varnothing$ 5 mm Al rod at 470 mm free length yields at \textbf{F\_y ${\approx}$ 7.20 N} (root yield moment M\_y = ${\sigma}$\_y ${\times}$ I/r = 276 ${\times}$ 10$^{6}$ ${\times}$ 3.068 ${\times}$ 10$^{-}$$^{1}$$^{1}$ / 2.5 ${\times}$ 10$^{-}$$^{3}$ = 3.386 N${\cdot}$m; F\_y = 3.386/0.470 = 7.20 N). Maximum safe test load is \textbf{5.0 N} (69 \% of yield). Do not exceed this.

\textbf{Procedure:}
\begin{enumerate}[noitemsep,topsep=2pt]
\item Clamp the root of one $\varnothing$ 5 mm, 500 mm long Al rod in a rigid bench vice with 30 mm grip overlap (free length L = 470 mm).
\item Apply a point load at the free tip using a precision hanging balance (suspended load in a cup via fishing line over a pulley aligned with the spar axis).
\item Record tip deflection vs. load in 0.5 N increments up to 5.0 N maximum:
\end{enumerate}

\begin{center}
\begin{tabular}{|l|l|l|}
\hline
Step & Load (N) & Expected tip deflection (mm) \\
\hline
1 & 1.0 & 16.4 \\
\hline
2 & 2.0 & 32.7 \\
\hline
3 & 3.0 & 49.1 \\
\hline
4 & 4.0 & 65.5 \\
\hline
5 & 5.0 & 81.8 - \textbf{Stop here} \\
\hline
\end{tabular}
\end{center}

Expected tip deflection (elastic, point load, cantilever, I = ${\pi}$r$^{4}$/4 = ${\pi}$(2.5 ${\times}$ 10$^{-}$$^{3}$)$^{4}$/4 = 3.068 ${\times}$ 10$^{-}$$^{1}$$^{1}$ m$^{4}$):

\[
\delta = \frac{FL^3}{3EI} = \frac{F \times (0.47)^3}{3 \times 69 \times 10^9 \times 3.068 \times 10^{-11}} = \frac{F \times 0.10393}{6.351 \times 10^{-3}} = 16.37\,F\ \text{mm}
\]

\begin{enumerate}[noitemsep,topsep=2pt]
\item Stop the test at 5.0 N. Do not increase load further.
\item Compare measured stiffness with theoretical: $k_{theory} = 3EI/L^3 = 6.351 \times 10^{-3}/0.10393 = 61.1\ \text{mN/mm}$.
\end{enumerate}

\textbf{Expected results:}

\begin{center}
\begin{tabularx}{\textwidth}{|l|l|X|}
\hline
Metric & Expected & Acceptance criterion \\
\hline
Tip stiffness $k$ & 61.1 mN/mm (= 61.1 N/m) & 51.9 - 70.3 mN/mm ($\pm$15 \%) \\
\hline
Load at onset of yield & $\approx$ 7.20 N & Verify by observing permanent set after 5.0 N test: residual $<$ 1 mm = elastic \\
\hline
Deflection linearity up to 4.0 N & $R^{2} > 0.99$ for $F$-$\delta$ plot & $R^{2} \geq 0.98$ \\
\hline
\end{tabularx}
\end{center}

\begin{figure}[H]
    \centering
    \includegraphics[width=0.905\linewidth]{Fig11—Phase-B2_Spar_Load–Deflection_Test.png}
    \caption{Phase-B2 expected spar load-deflection curve.}
    \label{fig:fig11}
\end{figure}

The 16.4 mm tip deflection per Newton for the $\varnothing$ 5 mm spar is 7.7${\times}$ smaller than the 126 mm/N of a $\varnothing$ 3 mm rod - consistent with the $(d_5/d_3)^4 = (5/3)^4 = 7.72$ ratio of second moments of area. The purpose of this test is to verify E through measured stiffness and confirm the rod batch meets material specification.

\textbf{Cautions:}
\begin{itemize}[noitemsep,topsep=2pt]
\item Do \textbf{not} exceed 5.0 N on the rod in this free-tip 470 mm configuration; the yield load is 7.20 N and permanent deformation will render the spar unusable.
\item Check that the fishing line is perpendicular to the spar at the load point - any angular error underestimates stiffness.
\item Do not kink the rod during clamping; inspect the free end under magnification for scratches or tool marks that could initiate fatigue cracks.
\end{itemize}

\subsection{Phase C - Full Wing Static Test}

\subsubsection{Test C1: Wing Tip Deflection Under Distributed Load}

The $\varnothing$ 5 mm front spar remains fully elastic at all design load levels. The root front-spar bending moment at design ultimate load (2.07 N${\cdot}$m) is only 61 \% of the yield moment (3.39 N${\cdot}$m), providing a yield safety margin of +64 \%. The full wing may therefore be safely loaded to 100 \% design limit load (6.13 N half-wing) and, if desired, to design ultimate load (9.20 N half-wing) during structural testing. Standard load increments (${\leq}$ 10 \% of ultimate per 30-second step) and general safety precautions from No.15.1 still apply.

The purpose of the test is to verify linear-elastic structural response of the assembled wing and measure tip deflection up to design limit loads; confirm structural integrity and no damage (debonds, cracks) under the full design load range.

\textbf{Procedure:}
\begin{enumerate}[noitemsep,topsep=2pt]
\item Mount the completed wing with the root joint fixed to a rigid wall bracket. The wing panel lies horizontally (spanwise) in the test rig, with the leading edge facing forward.
\item Apply loads via sand-bag blocks or calibrated weights placed at equally-spaced rib locations to simulate uniform distributed lift.
\item Total half-wing limit lift: $L_{1/2,\text{limit}} = 6.13\ \text{N}$. Divide equally among 11 rib bays (approximately 0.56 N per rib station for 100 \% limit).
\item Apply loads in four increments: 25 \%, 50 \%, 75 \%, and 100 \% of limit load. An additional single increment to 100 \% ultimate load (9.20 N) may be applied if all prior increments show linear response and no anomalies. At each increment:
\end{enumerate}
\begin{itemize}[noitemsep,topsep=2pt]
\item Measure tip deflection with a steel rule or dial gauge referenced to the root.
\item Photograph the wing from the root and tip directions.
\item Check for audible cracking (evidence of rib or skin delamination).
\end{itemize}
\begin{enumerate}[noitemsep,topsep=2pt]
\item Hold at 100 \% limit load for 60 seconds; record any creep.
\item Unload and measure permanent set (residual tip displacement).
\end{enumerate}

\textbf{Expected results (using I = 3.068 ${\times}$ 10$^{-}$$^{1}$$^{1}$ m$^{4}$ for $\varnothing$ 5 mm rod):}

\begin{center}
\begin{tabularx}{\textwidth}{|l|l|l|X|}
\hline
Load level & Applied load (N) & Expected tip deflection (mm) &  \\
\hline
25 \% limit & 1.53 & $\sim$11 & Linear; no visible deformation \\
\hline
50 \% limit & 3.07 & $\sim$23 & Linear \\
\hline
75 \% limit & 4.60 & $\sim$34 & Still elastic; continue if no anomalies observed \\
\hline
100 \% limit & 6.13 & $\sim$45 & Spar at 41 \% of yield moment; fully elastic \\
\hline
100 \% ultimate & 9.20 & $\sim$68 & Structural qualification load; spar at 61 \% yield \\
\hline
\end{tabularx}
\end{center}

\textbf{Calculation for reference:} Tip deflection under uniform distributed load at 100 \% limit load:
$$\delta\_{\text{tip}} = \frac{w\_0 (b/2)^4}{8EI} = \frac{12.26 \times (0.500)^4}{8 \times 69 \times 10^9 \times 3.068 \times 10^{-11}} = \frac{0.7663}{1.694 \times 10^{-2}} = 45.2\ \text{mm} \approx 9.0\%\ \text{semi-span}$$

The $\varnothing$ 5 mm front spar is structurally adequate at 1 m span - no yielding occurs at or below design ultimate load. The 9 \% semi-span deflection at limit load is large but elastically recoverable; it is characteristic of flexible solar-UAV-style construction and is accepted for this prototype.

\textbf{Pass/fail criteria (up to 100 \% limit):}
\begin{itemize}[noitemsep,topsep=2pt]
\item \textbf{Pass:} Deflection vs. load is linear (R$^{2}$ > 0.99), no audible cracking, residual permanent set < 1 mm after unloading, deflection at 100 \% limit is within ${\pm}$20 \% of predicted 45 mm.
\item \textbf{Fail:} Any audible crack, skin delamination, rib detachment, or permanent set > 2 mm after unloading to zero.
\end{itemize}

\textbf{Cautions:}
\begin{itemize}[noitemsep,topsep=2pt]
\item Ensure load application points are at rib centrelines, not between ribs, to avoid local skin dents.
\item Keep fingers and hands clear of the upper spar surface during loading.
\item Have a safety net or foam pad below the wing to catch falling weights if a fastener releases.
\end{itemize}

\begin{figure}[H]
    \centering
    \includegraphics[width=0.905\linewidth]{Fig12_Phase-C1_Wing_Static_Test_Deflection.png}
    \caption{Phase-C1 expected wing deflection shapes.}
    \label{fig:fig12}
\end{figure}

\subsubsection{Test C2: Wing Torsion Test}

The purpose is to verify the D-box torsional stiffness and confirm that the wing does not twist excessively under aerodynamic pitching moment.

\textbf{Procedure:}
\begin{enumerate}[noitemsep,topsep=2pt]
\item With the wing mounted root-fixed, apply a torque at 75 \% semi-span by hanging equal-and-opposite weights (nose-down/nose-up couple) at the leading edge and trailing edge of one rib, 300 mm from the root.
\item Target torque: 0.05 N${\cdot}$m (half the root torque calculated in No.10.1).
\item Measure twist angle using a digital inclinometer placed on the rib at the load station.
\item Increase torque in 10 \% steps to 0.10 N${\cdot}$m.
\end{enumerate}

\textbf{Expected results:}

\begin{center}
\begin{tabular}{|l|l|}
\hline
Torque (N${\cdot}$m) & Expected twist (degrees) \\
\hline
0.025 & ~0.007 \\
\hline
0.050 & ~0.014 \\
\hline
0.100 & ~0.029 \\
\hline
\end{tabular}
\end{center}

Analytical twist (from No.10.3): $\phi = 0.014^{\circ}$ for T\_root = 0.105 N${\cdot}$m. Values here are lower because load is applied at 75 \% span, not root.

\textbf{Acceptance criterion:} Twist per unit torque ${\leq}$ 0.3$^{\circ}$/N${\cdot}$m at the load station. Residual twist < 0.05$^{\circ}$ after load removal.

\textbf{Cautions:}
\begin{itemize}[noitemsep,topsep=2pt]
\item Ensure that the couple arm is rigid (use a stiff temporary balsa jig glued temporarily to the rib and removed after testing).
\item Do not apply torque in the trailing-edge-up (positive pitch) direction beyond 0.15 N${\cdot}$m, which risks peel-off of the D-box skin at the leading edge.
\end{itemize}

\subsubsection{Test C3: Tap Test - Skin/Rib Debond Detection}

The purpose is to detect any adhesive delamination or debonds between skin, ribs, and spars without destructive testing.

\textbf{Procedure:}
\begin{enumerate}[noitemsep,topsep=2pt]
\item After the static load tests (C1, C2), lightly tap each rib-to-skin joint and each rib-to-spar glue fillet using a small coin or pencil.
\item A solid "thud" indicates a well-bonded joint.
\item A hollow "click" or "pop" indicates a debond; mark and record location.
\end{enumerate}

\textbf{Expected results:} All joints should produce a solid thud. No more than 5 \% of joint area should show any hollow response after limit-load testing.

\textbf{Cautions:}
\begin{itemize}[noitemsep,topsep=2pt]
\item Apply tapping force of < 1 N; heavier taps can initiate fresh delaminations in the thin balsa.
\item Perform the tap test in a quiet environment; ambient noise masks the acoustic signature.
\end{itemize}

\subsection{Phase D - Vibration / Modal Test}

\subsubsection{Test D1: Free-Vibration Pluck Test (Wing Natural Frequency)}

The purpose is to confirm the first bending natural frequency is consistent with the analytical prediction (f$_{1}$ ${\approx}$ 5.97 Hz from No.12.1).

\textbf{Procedure:}
\begin{enumerate}[noitemsep,topsep=2pt]
\item Mount the wing root-fixed on the bench bracket.
\item "Pluck" the wing tip with a light tap (finger flick) and immediately release.
\item Record the free-decay vibration using a smartphone accelerometer app (or a small accelerometer and oscilloscope) attached near the wing tip.
\item Count zero-crossings or use an FFT to extract the dominant frequency.
\end{enumerate}

\textbf{Expected results:}

\begin{center}
\begin{tabular}{|l|l|l|}
\hline
Frequency & Expected & Acceptance criterion \\
\hline
1st bending f$_{1}$ & ~5.97 Hz & 4.5 - 7.5 Hz (${\pm}$25 \% for prototype variability) \\
\hline
Damping ratio $\zeta$ & 0.01 - 0.05 & < 0.10 (lightly damped) \\
\hline
\end{tabular}
\end{center}

\textbf{Cautions:}
\begin{itemize}[noitemsep,topsep=2pt]
\item Do not excite the wing with a resonant forcing frequency near f$_{1}$ for more than 2 seconds; resonance can amplify displacement to failure even at small forces.
\item Ensure no loose masses (weights, tools) are resting on the wing during the pluck test, as these shift the measured frequency.
\item If the measured frequency is below 4 Hz, inspect for debonded root joint or cracked spar - this indicates reduced stiffness and requires repair before any load testing.
\end{itemize}

\subsection{Test Record and Sign-Off}

All test results should be recorded in the test log with:
\begin{itemize}[noitemsep,topsep=2pt]
\item Date and tester name
\item Batch number/label of each material used
\item Load cell or weight calibration certificate reference
\item Photographs at each load step (minimum: unloaded, 25 \% limit, 50 \% limit, unloaded-after)
\item Pass/fail judgement for each acceptance criterion
\end{itemize}

A completed test log signed by the laboratory supervisor constitutes the airworthiness evidence for the prototype wing. No flight or wind-tunnel testing should proceed until all Phase A, B, and C tests have returned \textbf{Pass} results.

\newpage 

\appendix
\section{Appendix A - Material Properties Summary}

\begin{center}
\renewcommand{\arraystretch}{1.2}
\begin{tabularx}{\textwidth}{|X|p{1.5cm}|p{1.5cm}|p{2cm}|p{2cm}|p{2cm}|}
\hline
Material & $E$ (GPa) & $G$ (GPa) & $\sigma_y$ (MPa) & $\sigma_{ult}$ (MPa) & $\rho$ (kg/m$^{3}$) \\
\hline
Al 6061-T6 & 69 & 26 & 276 & 310 & 2700 \\
\hline
Balsa (along grain, medium) & 3.4 & 0.16 & 14.7 & 14.7 & 160 \\
\hline
Carbon fibre (unidirectional, for reference) & 135 & 5.1 & - & 1500 & 1600 \\
\hline
Polyester film (Oracover, for reference) & 3.5 & - & 50 & 70 & 1390 \\
\hline
\end{tabularx}
\end{center}

\newpage \section{Appendix B - Assumptions Correction Table}

\begin{center}
\renewcommand{\arraystretch}{1.2}
\begin{tabularx}{\textwidth}{|p{3cm}|p{3cm}|p{3.5cm}|X|}
\hline
Assumption & Initial (uncorrected) & Corrected value & Correction method \\
\hline
Total mass & 70 g (cubic scaling) & \textbf{500 g} & Structural weight model + component budget \\
\hline
Lift distribution & Uniform & \textbf{Elliptical} & Prandtl lifting-line (rectangular wing, $AR = 3.92$) \\
\hline
Root bending moment & 2.25 N$\cdot$m (uniform) & \textbf{1.91 N$\cdot$m} (elliptical) & 15\% reduction; use uniform as conservative design value \\
\hline
$C_P$ location & 25\% chord & \textbf{31.3\% chord} & From $C_{M,AC} = -0.05$, $C_L = 0.80$ \\
\hline
D-box area & Not estimated & \textbf{490.9 mm$^{2}$} & Semi-ellipse geometry of airfoil leading edge \\
\hline
Gust load factor & $n = 2.5$ & $\Delta n_{\text{gust}} = 7.31$, \textbf{limited to $n = 2.5$ operationally} & Calm-weather flight restriction \\
\hline
Skin in bending & Participates & \textbf{Spar only} (skin is secondary) & Comparison of skin/spar stiffness \\
\hline
Torsional stiffness & Open section & \textbf{Closed D-box} (Bredt-Batho) & Leading edge + front spar + skin form closed cell \\
\hline
\end{tabularx}
\end{center}

\newpage \section{Appendix C - Full-Scale Reference Analysis (Sky-Sailor)}

This appendix provides reference analysis of the \textbf{full-scale Sky-Sailor solar UAV} (3.2 m span, André Noth, ETH Zürich, 2008). It serves as a design benchmark and shows how the scale model relates to the full-scale design.


\subsection{Full-Scale Wing Parameters}


\begin{center}
\begin{tabular}{|l|l|l|l|}
\hline
Parameter & Symbol & Value & Unit \\
\hline
Wingspan & b & 3.200 & m \\
\hline
Semi-span & s & 1.600 & m \\
\hline
Chord & c & 0.250 & m \\
\hline
Wing area & S & 0.776 & m$^{2}$ \\
\hline
Aspect ratio & AR & 13.0 & - \\
\hline
MTOW & W & 2.444 ${\times}$ 9.81 = 23.96 & N \\
\hline
Wing loading & W/S & 31.59 & N/m$^{2}$ \\
\hline
Airfoil & - & WE 3.55-9.3 & - \\
\hline
t/c & - & 9.3 \% & - \\
\hline
Camber ratio & cam/c & 3.55 \% & - \\
\hline
Max thickness & t & 9.3 \% ${\times}$ 250 = 23.25 mm & - \\
\hline
Scale factor (vs prototype) & ${\lambda}$ & 0.490 / 1.600 = 0.306 & - \\
\hline
\end{tabular}
\end{center}

\textbf{Full-scale design load factor (after gust alleviation):}

Using FAR 23 No.23.337 / No.23.341 with K\_g:

\[
\mu_g = \frac{2(m/S)}{\rho\,c\,a} = \frac{2 \times (2.444/0.776)}{1.225 \times 0.250 \times 5.44} = \frac{6.30}{1.662} = 3.790
\]

\[
K_g = \frac{0.88 \times 3.790}{5.3 + 3.790} = \frac{3.335}{9.090} = 0.367
\]

\[
\Delta n_{\text{gust,raw}} = \frac{\rho\,V_c\,U_e\,a}{2(W/S)} = \frac{1.225 \times 8.3 \times 7.62 \times 5.44}{2 \times 31.59} = \frac{423.1}{63.18} = 6.70
\]

\[
\Delta n_{\text{gust,allev}} = 0.367 \times 6.70 = 2.46 \quad \Rightarrow \quad n_{\text{gust,total}} = 3.46 \approx n_{\text{manoeuvre}} = 3.5
\]

Design load factor: \textbf{n = 3.5} (FAR 23 manoeuvre limit governs after alleviation).


\subsection{Full-Scale Structural Loads}


\textbf{Inertia-relieved net load per semi-span (n = 3.5):}

\begin{align}
L_{\text{per wing}} &= \frac{n W}{2} 
= \frac{3.5 \times 23.96}{2} 
= 41.93\ \text{N} \\
\\
W_{\text{inertia}} &= m_{\text{wing, semi}} \, g \, n 
= 0.2625 \times 9.81 \times 3.5 
= 9.01\ \text{N} \\
\\
L_{\text{net}} &= L_{\text{per wing}} - W_{\text{inertia}} 
= 41.93 - 9.01 
= 32.92\ \text{N}
\end{align}

\textbf{Root values:}

\begin{center}
\begin{tabular}{|l|l|}
\hline
Quantity & Value \\
\hline
Root shear V(0) & 32.92 N \\
\hline
Root bending moment M(0) & 26.35 N${\cdot}$m \\
\hline
Root torque T(0) & -0.44 N${\cdot}$m (nose-down) \\
\hline
Root lift per unit span l$_{0}$ (elliptical) & 33.37 N/m \\
\hline
\end{tabular}
\end{center}


\subsection{Full-Scale Main Spar (CFRP I-Beam)}


The full-scale spar is a carbon-fibre-reinforced polymer (CFRP) I-section, replacing the aluminium rod used in the scale prototype.

Spar cross-section:
  Height h         = t\(_\text{max}\) = 23.25 mm  (cap-to-cap)
  Cap width b\(_\text{cap}\)  = 15 mm
  Cap thickness    = 1.0 mm (2 x 0.5 mm CFRP prepreg, T300)
  Web thickness    = 2.0 mm (balsa, grain vertical)
  Web height h\(_\text{web}\) = h - 2 x t\(_\text{cap}\) = 21.25 mm

\textbf{Second moment of area:}

\[
I_{\text{caps}} = 2 \times b_{\text{cap}} \times t_{\text{cap}} \times \left(\frac{h}{2} - \frac{t_{\text{cap}}}{2}\right)^2 = 2 \times 0.015 \times 0.001 \times (0.01163)^2 = 4.06 \times 10^{-9}\ \text{m}^4
\]

\[
I_{\text{web}} = \frac{t_{\text{web}} \times h_{\text{web}}^3}{12} = \frac{0.002 \times (0.02125)^3}{12} = 1.60 \times 10^{-9}\ \text{m}^4
\]

\[
I_{xx} = I_{\text{caps}} + I_{\text{web}} = 4.06 \times 10^{-9} + 1.60 \times 10^{-9} = 5.66 \times 10^{-9}\ \text{m}^4
\]

\textbf{Bending stress in CFRP caps (E\_CFRP = 120 GPa):}

\[
\sigma_{\max} = \frac{M_{\text{root}} \times y_{\max}}{I_{xx}} = \frac{26.35 \times 0.01163}{5.66 \times 10^{-9}} = \frac{0.3065}{5.66 \times 10^{-9}} = 54.2\ \text{MPa}
\]

Allowable (T300 CFRP UD): ${\sigma}$\_allow = 500 MPa (compression), giving \textbf{SF = 9.2} - caps are oversized from a stress standpoint; minimum practical cap size drives the design.

\textbf{Web shear stress:}

\[
\tau_{\text{web}} = \frac{V \times Q}{I_{xx} \times t_{\text{web}}} = \frac{32.92 \times 1.669 \times 10^{-7}}{5.66 \times 10^{-9} \times 0.002} = \frac{5.496 \times 10^{-6}}{1.132 \times 10^{-11}} = 0.485\ \text{MPa}
\]

Allowable shear (balsa, vertical grain): ${\tau}$\_allow = 3.0 MPa ${\rightarrow}$ \textbf{SF = 6.2} $\checkmark $

\textbf{Web shear buckling check (balsa web, 2.0 mm thick):}

\[
\tau_{\text{cr}} = \frac{k_s \pi^2 E_{\text{balsa}}}{12(1-\nu^2)} \left(\frac{t_{\text{web}}}{h_{\text{web}}}\right)^2 = \frac{5.35 \times 9.870 \times 3.5 \times 10^9}{10.92} \times \left(\frac{0.002}{0.02125}\right)^2
\]

\[
= 1.693 \times 10^{10} \times 8.873 \times 10^{-3} = 1.502 \times 10^8\ \text{Pa}... 
\]

This is the critical buckling stress for the web panel (Euler plate buckling with k\_s = 5.35 for infinitely long plate, a/b >> 1). \textbf{SF\_web\_buckling = 150 / 0.485 = 309} - web buckling is not a concern.

An earlier version of the spar used t\_web = 1.5 mm, which gave ${\tau}$\_cr = 84.4 MPa and SF\_buckling = 174. Even the thinner web is adequate. The web thickness of 2.0 mm was chosen for \textbf{ease of fabrication} (easier to cut and handle) rather than structural necessity.

\textbf{Full-scale spar structural weight:}


CFRP caps:  m = rho\(_{\text{CFRP}}\) x 2 x b\(_\text{cap}\) x t\(_\text{cap}\) x s = 1600 x 2 x 0.015 x 0.001 x 1.6 = 76.8 g\\
Balsa web:  m = rho\(_\text{balsa}\) x t\(_\text{web}\) x h\(_\text{web}\) x s     = 160 x 0.002 x 0.02125 x 1.6   = 10.9 g\\
Total main spar (per semi-span): 87.7 g\\

Comparison with Noth data: total half-wing (including skin, ribs, solar cells, servos) = 161.6 g. Spar mass = 87.7 g ${\approx}$ \textbf{54 \% of half-wing mass} - consistent with a spar-dominated construction.


\subsection{Full-Scale Rear Spar (CFRP Tube)}

Location: 65; chord = 162.5 mm from LE\\
Airfoil thickness at 65\(\%\) chord \(\approx\) 6\(\%\): t\(_\text{rear}\) = 0.06 x 250 = 15.0 mm\\
Section: CFRP rectangular tube 15 \(\times\) 6 mm, t\(_\text{wall}\) = 0.5 mm\\
\begin{align*}
I_{\text{rear}} &= \frac{6 \times 15^{3} - 5 \times 14^{3}}{12} \times 10^{-12} \\
&= \frac{6 \times 3375 - 5 \times 2744}{12} \times 10^{-12} \\
&= \frac{20250 - 13720}{12} \times 10^{-12} \\
&= \frac{6530}{12} \times 10^{-12} \\
&= 544.167 \times 10^{-12} \\
&= 5.442 \times 10^{-10}\ \text{m}^4
\end{align*}

Aileron hinge bending moment (20$^{\circ}$ deflection, aileron span = 30\% of semi-span = 0.48 m):

\[
M_{\text{RS,root}} = \Delta L_{\text{aileron}} \times \bar{y}_{\text{aileron}} = 4.19 \times 0.24 = 1.01\ \text{N}\cdot\text{m}
\]

\[
\sigma_{\text{rear}} = \frac{M \times y}{I} = \frac{1.01 \times 0.0075}{5.44 \times 10^{-10}} = 13.9\ \text{MPa} \ll 500\ \text{MPa}\ \checkmark
\]


\subsection{Full-Scale Natural Frequencies}

\begin{center}
\renewcommand{\arraystretch}{1.3}
\begin{tabularx}{\textwidth}{|p{2.5cm}|X|p{3cm}|}
\hline
Mode & Formula & Value \\
\hline
1st bending $f_{1}$ & $\dfrac{(1.875)^2}{2\pi L^2}\sqrt{\dfrac{EI}{\bar{m}}}$, with $EI = 120\times10^{9} \times 5.66\times10^{-9} = 679.2\ \text{N}\cdot\text{m}^2$, $\bar{m} = 0.1641\ \text{kg/m}$, $L = 1.6\ \text{m}$ & \textbf{7.6 Hz} \\
\hline
1st torsion $f_T$ & $\dfrac{\pi}{4L}\sqrt{\dfrac{GJ}{I_\alpha}}$, with $GJ = 12.22\ \text{N}\cdot\text{m}^2$, $I_\alpha = 8.62\times10^{-4}\ \text{kg}\cdot\text{m}$ & \textbf{18.6 Hz} \\
\hline
$f_T / f_B$ & $18.6 / 7.6 = 2.45$ & Flutter coupling not expected $\checkmark$ \\
\hline
\end{tabularx}
\end{center}

\subsection{Full-Scale Aeroelastic Margins}

\begin{center}
\renewcommand{\arraystretch}{1.2}
\begin{tabularx}{\textwidth}{|p{3cm}|p{3cm}|X|p{2.5cm}|}
\hline
Check & Value & Design requirement & Margin \\
\hline
Divergence speed $V_{\text{div}}$ & $\gg V_D$ (elastic axis $\approx$ AC) & $1.15 \times V_D = 13.4$ m/s & $\checkmark$ Infinite margin \\
\hline
Flutter speed $V_{\text{flutter}}$ & $\sim 47$ m/s (simplified estimate) & $1.20 \times V_D = 13.9$ m/s & $\checkmark\ \times 3.4$ \\
\hline
Tip twist at $n = 3.5$ & $3.30^{\circ}$ & $< 5^{\circ}$ & $\checkmark$ \\
\hline
Tip deflection at $n = 3.5$ & 22.9 mm (1.43 \% of semi-span) & $< 5\%$ semi-span & $\checkmark$ \\
\hline
\end{tabularx}
\end{center}

\subsection{Comparison: Full-Scale vs Scale Prototype}

\begin{center}
\begin{tabular}{|l|l|l|}
\hline
Parameter & Full-Scale (Sky-Sailor) & Scale Prototype (Group 2) \\
\hline
Span & 3.20 m & 0.98 m (${\lambda}$ = 0.306) \\
\hline
Chord & 0.250 m & 0.250 m \\
\hline
MTOW & 23.96 N & 4.905 N \\
\hline
Design load factor n & 3.5 & 2.5 \\
\hline
Airfoil & WE3.55-9.3 (cambered) & WE3.55-9.3 \\
\hline
Main spar & CFRP I-beam, h = 23.25 mm & Al 6061-T6 rod, $\varnothing$ 5 mm \\
\hline
Root bending moment & 26.35 N${\cdot}$m & 2.025 N${\cdot}$m \\
\hline
Root shear & 32.92 N & 9.20 N \\
\hline
Spar bending SF & 8.03 & 1.52 ($\varnothing$ 5 mm) \\
\hline
1st bending frequency & 7.6 Hz & 5.97 Hz \\
\hline
1st torsion frequency & 18.6 Hz & 145.8 Hz \\
\hline
Tip deflection & 22.9 mm (1.4 \% s) & ~16.9 mm (3.4 \% s) at limit \\
\hline
Wing mass (per half) & ~101.5 g (full structure) & ~50 g (estimated) \\
\hline
\end{tabular}
\end{center}

\newpage

\begin{thebibliography}{99}

\bibitem{Noth2008}
A. Noth and R. Siegwart, \textit{Design of Solar Powered Airplanes for Continuous Flight}, Dissertation No. 18010, ETH Zürich, Zürich, Switzerland, 2008.

\bibitem{Young2002}
W. C. Young and R. G. Budynas, \textit{Roark's Formulas for Stress and Strain}, 7th ed., McGraw-Hill, New York, NY, 2002.

\bibitem{Megson4thed}
T. H. G. Megson, \textit{Aircraft Structures for Engineering Students}, 4th ed., Butterworth-Heinemann, Oxford, UK.

\bibitem{Maughmer}
M. D. Maughmer and S. A. Brandt, \textit{Aerodynamics for Engineers}, Prentice Hall, Upper Saddle River, NJ.

\bibitem{Blevins1979}
R. D. Blevins, \textit{Formulas for Natural Frequency and Mode Shape}, Van Nostrand Reinhold, New York, NY, 1979.

\bibitem{FAR23}
\textit{Federal Aviation Regulations}, Part 23, Sections 23.335, 23.337, 23.341, Federal Aviation Administration, Washington, DC.

\bibitem{Temoshenko1961}
S. Timoshenko and J. M. Gere, \textit{Theory of Elastic Stability}, 2nd ed., McGraw-Hill, New York, NY, 1961.

\bibitem{Prandtl1918}
L. Prandtl, “Tragflügeltheorie (Lifting-Line Theory),” Nachrichten der Gesellschaft der Wissenschaften zu Göttingen, 1918.

\bibitem{ICAODoc7488}
\textit{Manual of the ICAO Standard Atmosphere}, 3rd ed., ICAO Doc 7488, International Civil Aviation Organization, Montreal, Canada, 1993.

\end{thebibliography}

\end{document}
